% EB presentation, 13 November 2026. Build: bash build.sh
% Every result number is a macro from numbers.tex (scripts/eb_numbers.py).
\def\talkbuild{1}
\documentclass[xcolor=table,aspectratio=169]{beamer}

\usepackage[defaultsans,scale=0.95]{opensans}
\usepackage{FiraMono}
\usepackage{fourier}
\usefonttheme{professionalfonts}
\DeclareMathAlphabet{\mathcal}{OMS}{cmsy}{m}{n}
\newcommand\bmmax{2}
\usepackage{bm}
\usepackage[T1]{fontenc}

\usepackage{graphicx}
\usepackage{caption}
\captionsetup[figure]{labelformat=empty}
\usepackage{subcaption}
\usepackage{multicol}

\usepackage{tikz}
\usetikzlibrary{decorations.pathreplacing}
\usepackage{wrapfig}

% COLORS
\usepackage{xcolor}
\colorlet{myred}{red!80!black}
\colorlet{myblue}{blue!80!black}
\colorlet{mygreen}{green!60!black}
\colorlet{myorange}{orange!70!red!60!black}
\colorlet{mydarkred}{red!30!black}
\colorlet{mydarkblue}{blue!40!black}
\colorlet{mydarkgreen}{green!30!black}
\colorlet{mylightblue}{blue!60!white}
\colorlet{mydarkyellow}{yellow!70!black}
\colorlet{myyellowred}{yellow!70!red}

\usepackage{float}
\usepackage{hyperref}

\usepackage{amsmath, amsthm, amssymb, amsfonts}
\usepackage{esint}
\usepackage{mathtools}
\usepackage{listings}
\lstset{basicstyle=\ttfamily}

\usepackage{physics}
\usepackage[makeroom]{cancel}
% siunitx + mhchem commented out due to version clash with chemgreek on this TL2017 install
\usepackage[italic]{heppennames}
\usepackage{xpatch}
\makeatletter
\xpatchcmd\@HepConStyle
{\edef\@upcode{\updefault}}
{\ifdefined\shapedefault\edef\@upcode{\shapedefault}\else\edef\@upcode{\updefault}\fi}
{}{}
\makeatother
\usepackage{tensor}
\usepackage{feynmp-auto}
\usepackage{listofitems}
\usetikzlibrary{arrows.meta}
\usepackage[outline]{contour}
\usepackage{booktabs}

%%% ---- shared TikZ vocabulary for the Part I method figures ----
\usetikzlibrary{positioning,calc,backgrounds}
\tikzset{
    pvflow/.style   ={-{Latex[length=2.4mm,width=1.9mm]}, line width=0.9pt, draw=black!55},
    pvpoint/.style  ={-{Latex[length=2.0mm,width=1.6mm]}, line width=0.9pt},
    pvaxis/.style   ={-{Latex[length=1.9mm,width=1.5mm]}, line width=0.7pt, draw=black!60},
    pvdim/.style    ={{Latex[length=1.7mm,width=1.4mm]}-{Latex[length=1.7mm,width=1.4mm]},
                      line width=0.7pt},
    pvcurve/.style  ={line width=1.0pt, line join=round, line cap=round},
    pvthresh/.style ={dash pattern=on 2.6pt off 2.2pt, line width=0.8pt,
                      draw=myorange!85!black},
    pvsplit/.style  ={dash pattern=on 3pt off 2.4pt, line width=1.0pt, draw=pennred},
    pvguide/.style  ={line width=0.5pt, draw=black!15},
    pvcard/.style   ={rounded corners=6pt, draw=black!15, fill=black!2, line width=0.7pt},
    pvbox/.style    ={rounded corners=6pt, align=center, line width=0.9pt},
}
\newcommand{\pvcap}[1]{{\scriptsize\textcolor{black!55}{#1}}}


\addtobeamertemplate{frametitle}{}{\vspace{-15.5pt}}

%%%%%%%%%%%%%%%%%%%%%%%%%%%%%%%%%%%%%%

\newcommand{\backupbegin}{
    \newcounter{finalframe}
    \setcounter{finalframe}{\value{framenumber}}
    \setbeamertemplate{footline}{}
    \begin{frame}[plain]
        \centering
        \vfill
        {\color{pennred!28}\rule{0.30\linewidth}{1.1pt}}\\[0.8em]
        {\LARGE\textbf{\textcolor{pennred}{Backup}}}\\[0.8em]
        {\color{pennred!28}\rule{0.30\linewidth}{1.1pt}}
        \vfill
    \end{frame}
}
\newcommand{\backupend}{
    \setcounter{framenumber}{\value{finalframe}}
    \setbeamertemplate{footline}[frame number]
}

\newcommand{\ntrk}{\ensuremath{n_\mathrm{trk}}}

\def\thisframelogos{}
\newcommand{\framelogo}[1]{\def\thisframelogos{#1}}

\addtobeamertemplate{frametitle}{}{%
\begin{tikzpicture}[remember picture,overlay]
\node[anchor=north east,yshift=3.5pt] at (current page.north east) {%
    \foreach \img in \thisframelogos {%
        \hspace{.5ex}%
        \includegraphics[height=0.8cm, keepaspectratio]{\img}%
    }%
};
\end{tikzpicture}}

% Asset folders are symlinked locally (e.g. logo-ATLAS -> ../../legacy/...)

% ---------------------------------------------------------------------------
% Two builds from one source.  See build.sh.
%
%   slides.pdf       reading version   (default, \talkbuild = 0)
%   slides_talk.pdf  presentation      (pdflatex "\def\talkbuild{1}% EB presentation, 13 November 2026. Build: bash build.sh
% Every result number is a macro from numbers.tex (scripts/eb_numbers.py).
\def\talkbuild{1}
\documentclass[xcolor=table,aspectratio=169]{beamer}

\usepackage[defaultsans,scale=0.95]{opensans}
\usepackage{FiraMono}
\usepackage{fourier}
\usefonttheme{professionalfonts}
\DeclareMathAlphabet{\mathcal}{OMS}{cmsy}{m}{n}
\newcommand\bmmax{2}
\usepackage{bm}
\usepackage[T1]{fontenc}

\usepackage{graphicx}
\usepackage{caption}
\captionsetup[figure]{labelformat=empty}
\usepackage{subcaption}
\usepackage{multicol}

\usepackage{tikz}
\usetikzlibrary{decorations.pathreplacing}
\usepackage{wrapfig}

% COLORS
\usepackage{xcolor}
\colorlet{myred}{red!80!black}
\colorlet{myblue}{blue!80!black}
\colorlet{mygreen}{green!60!black}
\colorlet{myorange}{orange!70!red!60!black}
\colorlet{mydarkred}{red!30!black}
\colorlet{mydarkblue}{blue!40!black}
\colorlet{mydarkgreen}{green!30!black}
\colorlet{mylightblue}{blue!60!white}
\colorlet{mydarkyellow}{yellow!70!black}
\colorlet{myyellowred}{yellow!70!red}

\usepackage{float}
\usepackage{hyperref}

\usepackage{amsmath, amsthm, amssymb, amsfonts}
\usepackage{esint}
\usepackage{mathtools}
\usepackage{listings}
\lstset{basicstyle=\ttfamily}

\usepackage{physics}
\usepackage[makeroom]{cancel}
% siunitx + mhchem commented out due to version clash with chemgreek on this TL2017 install
\usepackage[italic]{heppennames}
\usepackage{xpatch}
\makeatletter
\xpatchcmd\@HepConStyle
{\edef\@upcode{\updefault}}
{\ifdefined\shapedefault\edef\@upcode{\shapedefault}\else\edef\@upcode{\updefault}\fi}
{}{}
\makeatother
\usepackage{tensor}
\usepackage{feynmp-auto}
\usepackage{listofitems}
\usetikzlibrary{arrows.meta}
\usepackage[outline]{contour}
\usepackage{booktabs}

%%% ---- shared TikZ vocabulary for the Part I method figures ----
\usetikzlibrary{positioning,calc,backgrounds}
\tikzset{
    pvflow/.style   ={-{Latex[length=2.4mm,width=1.9mm]}, line width=0.9pt, draw=black!55},
    pvpoint/.style  ={-{Latex[length=2.0mm,width=1.6mm]}, line width=0.9pt},
    pvaxis/.style   ={-{Latex[length=1.9mm,width=1.5mm]}, line width=0.7pt, draw=black!60},
    pvdim/.style    ={{Latex[length=1.7mm,width=1.4mm]}-{Latex[length=1.7mm,width=1.4mm]},
                      line width=0.7pt},
    pvcurve/.style  ={line width=1.0pt, line join=round, line cap=round},
    pvthresh/.style ={dash pattern=on 2.6pt off 2.2pt, line width=0.8pt,
                      draw=myorange!85!black},
    pvsplit/.style  ={dash pattern=on 3pt off 2.4pt, line width=1.0pt, draw=pennred},
    pvguide/.style  ={line width=0.5pt, draw=black!15},
    pvcard/.style   ={rounded corners=6pt, draw=black!15, fill=black!2, line width=0.7pt},
    pvbox/.style    ={rounded corners=6pt, align=center, line width=0.9pt},
}
\newcommand{\pvcap}[1]{{\scriptsize\textcolor{black!55}{#1}}}


\addtobeamertemplate{frametitle}{}{\vspace{-15.5pt}}

%%%%%%%%%%%%%%%%%%%%%%%%%%%%%%%%%%%%%%

\newcommand{\backupbegin}{
    \newcounter{finalframe}
    \setcounter{finalframe}{\value{framenumber}}
    \setbeamertemplate{footline}{}
    \begin{frame}[plain]
        \centering
        \vfill
        {\color{pennred!28}\rule{0.30\linewidth}{1.1pt}}\\[0.8em]
        {\LARGE\textbf{\textcolor{pennred}{Backup}}}\\[0.8em]
        {\color{pennred!28}\rule{0.30\linewidth}{1.1pt}}
        \vfill
    \end{frame}
}
\newcommand{\backupend}{
    \setcounter{framenumber}{\value{finalframe}}
    \setbeamertemplate{footline}[frame number]
}

\newcommand{\ntrk}{\ensuremath{n_\mathrm{trk}}}

\def\thisframelogos{}
\newcommand{\framelogo}[1]{\def\thisframelogos{#1}}

\addtobeamertemplate{frametitle}{}{%
\begin{tikzpicture}[remember picture,overlay]
\node[anchor=north east,yshift=3.5pt] at (current page.north east) {%
    \foreach \img in \thisframelogos {%
        \hspace{.5ex}%
        \includegraphics[height=0.8cm, keepaspectratio]{\img}%
    }%
};
\end{tikzpicture}}

% Asset folders are symlinked locally (e.g. logo-ATLAS -> ../../legacy/...)

% ---------------------------------------------------------------------------
% Two builds from one source.  See build.sh.
%
%   slides.pdf       reading version   (default, \talkbuild = 0)
%   slides_talk.pdf  presentation      (pdflatex "\def\talkbuild{1}\input{slides}")
%
% \readingonly{...} keeps its argument only in the reading build; \talkonly{...}
% only in the talk build.  Every number lives in exactly one place, so the two
% versions cannot drift apart.
\providecommand{\talkbuild}{0}
\newcommand{\readingonly}[1]{\ifnum\talkbuild=0\relax#1\fi}
\newcommand{\talkonly}[1]{\ifnum\talkbuild=1\relax#1\fi}
% \bnum{reading}{talk} expands to a bare number, so it can be used inside a
% length: \includegraphics[height=\bnum{0.5}{0.6}\textheight]{...}.  The talk
% build drops prose, which frees vertical room for the figure that is left.
\newcommand{\bnum}[2]{\ifnum\talkbuild=0 #1\else #2\fi}
% ---------------------------------------------------------------------------

% Part-divider helper: \partdivider{numeral}{title}{one-line subtitle}
\newcommand{\partdivider}[3]{%
    \begin{frame}[plain]
        \centering
        \vfill
        \begin{tikzpicture}
            \fill[pennred] (0,0) circle (0.72cm);
            \node[text=white, font=\Large\bfseries] at (0,0.02) {#1};
        \end{tikzpicture}\\[0.85em]
        {\LARGE\textbf{\textcolor{pennred}{#2}}}\\[0.75em]
        {\color{pennred!28}\rule{0.30\linewidth}{1.1pt}}\\[0.75em]
        {\small\textcolor{black!50}{#3}}
        \vfill
    \end{frame}%
}

%%%%%%%%%%%%%%%%%%%%%%%%%%%%%%%%%%%%%%

\usetheme{Philadelphia}
\hypersetup{colorlinks,allcolors=pennred}

\title[PV-Finder]{PV-Finder: primary vertex reconstruction and track association with neural networks}
\subtitle{ANA-IDTR-2026-01, Editorial Board}

\author[Mattia Mauro]{Qi-Bin Lei, Rocky B. Garg, Lauren Tompkins,\\\textcolor{pennred}{Mattia Mauro}, Trey Warner}

\institute[]{Stanford University}

\date{November 13, 2026}

\titlegraphic{
    \hspace{5pt}
    \href{https://stanford.edu/group/stanford_atlas/}{\includegraphics[height=1cm]{logo-Stanford/Block_S_2_color.png}}
    \href{https://atlas.cern/}{\includegraphics[trim={0 11.5cm 0 11.5cm},clip,height=1cm]{logo-ATLAS/ATLAS-Logo-Ref-Pantone.pdf}}
    \href{https://iris-hep.org/}{\includegraphics[height=1cm]{logo-iris-hep/Iris-hep-4-no-long-name.png}}
}


\newcommand{\readingnote}[1]{\readingonly{\par\vspace{.35em}%
  {\footnotesize\color{black!65}#1\par}}}
\newcommand{\takeaway}[1]{\par\vspace{.65em}{\centering
  \setlength{\fboxsep}{7pt}\colorbox{pennred!5}{\parbox{.94\linewidth}{%
  \centering\small\textcolor{pennred}{#1}}}\par}}
\newcommand{\context}[1]{{\centering\footnotesize\color{black!55}#1\par}\vspace{.65em}}
\newcommand{\bigvalue}[2]{{\LARGE\bfseries\textcolor{pennred}{#1}}\\[.4em]{\small #2}\par}
\newcommand{\pubnote}{\href{https://cds.cern.ch/record/2858348}{ATL-PHYS-PUB-2023-011}}
\newcommand{\notelink}{\href{https://gitlab.cern.ch/atlas-physics-office/IDTR/ANA-IDTR-2026-01/ANA-IDTR-2026-01-INT1}{ANA-IDTR-2026-01-INT1}}
\newcommand{\treysource}{\href{https://gitlab.cern.ch/atlas-physics-office/IDTR/ANA-IDTR-2026-01/ANA-IDTR-2026-01-INT1/-/blob/4c50ed51500c897dd6a1b5bb52b0f1e3c12e7575/sections/04b_timing.tex}{Note Sec. 5.1 (18 Sep 2026)}}

\newcommand{\svtx}{\sigma_{\mathrm{vtx-vtx}}}
\setbeamertemplate{navigation symbols}{}
\hypersetup{pdftitle={PV-Finder, EB presentation, November 13, 2026},
  pdfauthor={Mattia Mauro, Qi-Bin Lei, Rocky B. Garg, Lauren Tompkins, Trey Warner}}

% numbers.tex: GENERATED by scripts/eb_numbers.py. Do not edit by hand.
%
% [density] checkpoint: hllhc_alleta_v6_mse_2ep_phase2_epoch_2 (v6 production)
%     window: 0.1694 mm (to become 0.17 mm)
%     peaks: amplitude 0.01, width >= 3 bins, integral 0.30, height 0.03, centroid
% [gnn] checkpoint: ttva_gnn_hllhc_v4_noaug epoch 107 (to be replaced by full-mu v2)
%     t: t098
% [lowpu] checkpoint: Run 2 reproduction model (reproduction_ep300), integral 0.50
%
% density: outputs/11_09_2026_output/split_inclusive_efficiency/results.json :: windows/r16443@0.1694/pvf/eff
\newcommand{\effPVF}{83.2}
% density: outputs/11_09_2026_output/split_inclusive_efficiency/results.json :: windows/r16443@0.1694/amvf/eff
\newcommand{\effAMVF}{79.8}
% density: outputs/11_09_2026_output/split_inclusive_efficiency/results.json :: windows/r16443@0.1694/diff/eff
\newcommand{\effDiff}{3.36}
% density: outputs/11_09_2026_output/split_inclusive_efficiency/results.json :: windows/r16443@0.1694/diff/eff_err
\newcommand{\effDiffErr}{0.06}
% density: outputs/11_09_2026_output/split_inclusive_efficiency/results.json :: windows/r16443@0.1/pvf/eff
\newcommand{\effHSwPVF}{74.1}
% density: outputs/11_09_2026_output/split_inclusive_efficiency/results.json :: windows/r16443@0.1/amvf/eff
\newcommand{\effHSwAMVF}{71.0}
% density: outputs/11_09_2026_output/split_inclusive_efficiency/results.json :: windows/r16443@0.1/diff/eff
\newcommand{\effHSwDiff}{3.12}
% density: outputs/11_09_2026_output/split_inclusive_efficiency/results.json :: windows/r16443@0.1694/pvf/fake_per_evt
\newcommand{\fakePVF}{18.7}
% density: outputs/11_09_2026_output/split_inclusive_efficiency/results.json :: windows/r16443@0.1694/amvf/fake_per_evt
\newcommand{\fakeAMVF}{19.7}
% density: outputs/11_09_2026_output/split_inclusive_efficiency/results.json :: windows/r16443@0.1694/pvf/split_per_evt
\newcommand{\splitPVF}{0.53}
% density: outputs/11_09_2026_output/split_inclusive_efficiency/results.json :: windows/r16443@0.1694/amvf/split_per_evt
\newcommand{\splitAMVF}{0.18}
% density: outputs/11_09_2026_output/split_inclusive_efficiency/results.json :: windows/r16443@0.1694/pvf/truth_per_evt
\newcommand{\truthPerEvt}{111.3}
% density: outputs/11_09_2026_output/split_inclusive_efficiency/results.json :: windows/r16443@0.1694/pvf/reco_per_evt
\newcommand{\recoPVF}{101.0}
% density: outputs/11_09_2026_output/split_inclusive_efficiency/results.json :: windows/r16443@0.1694/amvf/reco_per_evt
\newcommand{\recoAMVF}{97.9}
% density: outputs/11_09_2026_output/split_inclusive_efficiency/results.json :: windows/r16443@0.1694/n_events
\newcommand{\nEvtPU}{1{,}920}
% density: outputs/11_09_2026_output/split_inclusive_efficiency/results.json :: windows/r16443@0.1694/totals/pvf/n_truth
\newcommand{\nTruthPU}{213{,}738}
% density: outputs/11_09_2026_output/run2_vs_run4_resolution/results.json :: sigma/hl_pu200/pvf_int030_h003_cen/sigma_mm (integral 0.30, same as counting (one threshold set))
\newcommand{\sigPVF}{0.218}
% density: outputs/11_09_2026_output/run2_vs_run4_resolution/results.json :: sigma/hl_pu200/pvf_int030_h003_cen/fit_err_mm
\newcommand{\sigPVFErr}{0.003}
% density: outputs/11_09_2026_output/run2_vs_run4_resolution/results.json :: sigma/hl_pu200/amvf/sigma_mm
\newcommand{\sigAMVF}{0.305}
% density: outputs/11_09_2026_output/run2_vs_run4_resolution/results.json :: sigma/hl_pu200/amvf/fit_err_mm
\newcommand{\sigAMVFErr}{0.006}
% gnn: outputs/11_09_2026_output/hs_chain_study/hs_chain_summary.json :: reco_min_tracks_1/pu200/fractions/a_reco_pvf_peaks/value
\newcommand{\hsRecoPVF}{99.9}
% gnn: outputs/11_09_2026_output/hs_chain_study/hs_chain_summary.json :: reco_min_tracks_1/pu200/fractions/a_reco_amvf/value
\newcommand{\hsRecoAMVF}{99.8}
% gnn: outputs/11_09_2026_output/hs_chain_study/hs_chain_summary.json :: reco_min_tracks_1/pu200/mean_mu
\newcommand{\meanMu}{192.5}
% density: outputs/11_10_2026_output/small_model_eval/compare3_r16443.json :: s5/d_eff_pt
\newcommand{\sFiveDEffA}{-0.33}
% density: outputs/11_10_2026_output/small_model_eval/compare3_r16638.json :: s5/d_eff_pt
\newcommand{\sFiveDEffB}{-0.26}
% density: outputs/11_10_2026_output/small_model_eval/compare3_r16443.json :: s5/d_fp
\newcommand{\sFiveDFake}{-0.9}
% density: outputs/11_10_2026_output/small_model_eval/compare3_r16443.json :: s8/d_eff_pt
\newcommand{\sEightDEffA}{-0.80}
% density: outputs/11_10_2026_output/small_model_eval/compare3_r16638.json :: s8/d_eff_pt
\newcommand{\sEightDEffB}{-0.70}
% density: outputs/11_10_2026_output/small_model_eval/compare3_r16443.json :: s8/d_fp
\newcommand{\sEightDFake}{-3.1}
% density: outputs/11_09_2026_output/cpu_bench_packed/cpu_bench.json :: 0/median_ms
\newcommand{\cpuVsixPad}{1{,}601}
% density: outputs/11_09_2026_output/cpu_bench_packed/cpu_bench.json :: 1/median_ms
\newcommand{\cpuVsixPack}{1{,}159}
% density: outputs/11_09_2026_output/cpu_bench_packed/cpu_bench.json :: 3/median_ms
\newcommand{\cpuSFivePack}{250}
% density: outputs/11_09_2026_output/cpu_bench_packed/cpu_bench.json :: 5/median_ms
\newcommand{\cpuSEightPack}{178}
% gnn: outputs/11_09_2026_output/hs_chain_study/hs_chain_summary.json :: reco_min_tracks_1/pu200/clean_per_truth/gnn_v4_t098/value
\newcommand{\cleanTruthGNN}{68.4}
% gnn: outputs/11_09_2026_output/hs_chain_study/hs_chain_summary.json :: reco_min_tracks_1/pu200/clean_per_truth/amvf/value
\newcommand{\cleanTruthAMVF}{30.9}
% gnn: outputs/11_09_2026_output/hs_chain_study/hs_chain_summary.json :: reco_min_tracks_1/pu200/clean_per_truth/z0_nearest/value
\newcommand{\cleanTruthZnear}{31.9}
% gnn: outputs/11_09_2026_output/hs_chain_study/hs_chain_summary.json :: reco_min_tracks_1/pu200/clean_per_truth/z0_3sigma/value
\newcommand{\cleanTruthZsig}{33.1}
% gnn: outputs/11_09_2026_output/hs_chain_study/hs_chain_summary.json :: reco_min_tracks_1/pu200/fractions/b_clean_hs_gnn_v4_t098/value
\newcommand{\cleanHSGNN}{96.4}
% gnn: outputs/11_09_2026_output/hs_chain_study/hs_chain_summary.json :: reco_min_tracks_1/pu200/fractions/b_clean_hs_amvf/value
\newcommand{\cleanHSAMVF}{92.5}
% gnn: outputs/11_09_2026_output/hs_chain_study/hs_chain_summary.json :: reco_min_tracks_1/pu200/c_hs_track_fraction/gnn_v4_t098/value
\newcommand{\hsTrkGNN}{74.9}
% gnn: outputs/11_09_2026_output/hs_chain_study/hs_chain_summary.json :: reco_min_tracks_1/pu200/c_hs_track_fraction/amvf/value
\newcommand{\hsTrkAMVF}{74.1}
% gnn: outputs/11_09_2026_output/hs_chain_study/hs_chain_summary.json :: reco_min_tracks_1/pu200/c_hs_vertex_purity/gnn_v4_t098/value
\newcommand{\hsPurGNN}{91.3}
% gnn: outputs/11_09_2026_output/hs_chain_study/hs_chain_summary.json :: reco_min_tracks_1/pu200/c_hs_vertex_purity/amvf/value
\newcommand{\hsPurAMVF}{88.4}
% gnn: outputs/11_09_2026_output/hs_chain_study/hs_chain_summary.json :: reco_min_tracks_1/pu200/fractions/d_select_gnn_v4_t098_itk_0p1/value
\newcommand{\hsSelGNN}{96.3}
% gnn: outputs/11_09_2026_output/hs_chain_study/hs_chain_summary.json :: reco_min_tracks_1/pu200/fractions/d_native_amvf_itk_0p1/value (AMVF stored primary, same code as the GNN number)
\newcommand{\hsSelAMVF}{92.9}
% gnn: outputs/11_09_2026_output/hs_selection_reconciliation/reconciliation.json :: pu200_185_200/flags/rerank_plurality/value
\newcommand{\hsSelAMVFrerank}{80.7}
% gnn: outputs/11_11_2026_output/hs_position_estimators/results.json :: estimators/ada3d_t098/rms_um (3D adaptive fit with beam spot on the GNN tracks at the same t)
\newcommand{\hsResGNN}{9.7}
% gnn: outputs/11_11_2026_output/hs_position_estimators/results.json :: estimators/ada3d_t098/boot_err_um
\newcommand{\hsResGNNErr}{0.3}
% gnn: outputs/11_11_2026_output/hs_position_estimators/results.json :: estimators/amvf/rms_um
\newcommand{\hsResAMVF}{10.2}
% gnn: outputs/11_11_2026_output/hs_position_estimators/results.json :: estimators/amvf/boot_err_um
\newcommand{\hsResAMVFErr}{0.3}
% gnn: outputs/11_11_2026_output/hs_position_estimators/results.json :: estimators/centroid/rms_um
\newcommand{\hsResPeak}{12.4}
% gnn: outputs/11_11_2026_output/hs_position_estimators/results.json :: estimators/ada3d_truth/rms_um
\newcommand{\hsResTruth}{9.1}
% gnn: outputs/11_11_2026_output/hs_position_estimators/results.json :: n_common
\newcommand{\hsResNcommon}{1{,}915}
% lowpu: outputs/11_09_2026_output/data_mc_validation/results.json :: sigma/mc/fit
\newcommand{\sigMC}{0.416}
% lowpu: outputs/11_09_2026_output/data_mc_validation/results.json :: sigma/run2/fit
\newcommand{\sigRunTwo}{0.368}
% lowpu: outputs/11_09_2026_output/data_mc_validation/results.json :: sigma/run3/fit
\newcommand{\sigRunThree}{0.348}
% lowpu: outputs/11_09_2026_output/data_mc_validation/results.json :: matching/fitted_sample/mc_truth_ref/recovery
\newcommand{\effMC}{91.6}
% lowpu: outputs/11_09_2026_output/data_mc_validation/results.json :: matching/fitted_sample/mc_truth_ref/no_ref_per_evt
\newcommand{\noRefMCtruth}{0.42}
% lowpu: outputs/11_09_2026_output/data_mc_validation/results.json :: matching/fitted_sample/mc_amvf_ref/no_ref_per_evt
\newcommand{\noRefMCamvf}{1.72}
% lowpu: outputs/11_09_2026_output/data_mc_validation/results.json :: matching/fitted_sample/run2_data_amvf_ref/no_ref_per_evt
\newcommand{\noRefRunTwo}{1.82}
% lowpu: outputs/11_09_2026_output/data_mc_validation/results.json :: matching/fitted_sample/run3_data_amvf_ref/no_ref_per_evt
\newcommand{\noRefRunThree}{1.96}
% lowpu: outputs/11_09_2026_output/run2_vs_run4_resolution/results.json :: sigma/run2/amvf/sigma_mm
\newcommand{\sigMCAMVF}{0.789}
% lowpu: outputs/11_09_2026_output/run2_vs_run4_resolution/results.json :: sigma/run2/pvf_int050_h000_full/sigma_mm
\newcommand{\sigMCPVFr}{0.410}
% lowpu: outputs/11_09_2026_output/run2_vs_run4_resolution/results.json :: sigma/run2/ratio_amvf_over_pvf/ratio
\newcommand{\ratioRunTwo}{1.92}
% lowpu: outputs/11_11_2026_output/target_readout/results.json :: run2/target_run2cfg/sigma/sigma_mm
\newcommand{\tgtSigRunTwo}{0.203}
% lowpu: outputs/11_11_2026_output/target_readout/results.json :: hl_pu200/target_notecfg/sigma/sigma_mm
\newcommand{\tgtSigHL}{0.145}
% lowpu: outputs/11_11_2026_output/target_readout/results.json :: hl_pu200/target_notecfg_w1/sigma/sigma_mm
\newcommand{\tgtSigHLwone}{0.137}
% lowpu: outputs/11_11_2026_output/target_readout/results.json :: hl_pu200/target_notecfg_w1_g002/sigma/sigma_mm
\newcommand{\tgtSigHLfine}{0.124}
% lowpu: outputs/11_11_2026_output/target_readout/results.json :: ideal_precision_um/run2/pu_median
\newcommand{\idealPURunTwo}{75.6}
% lowpu: outputs/11_11_2026_output/target_readout/results.json :: ideal_precision_um/hl/pu_median
\newcommand{\idealPUHL}{44.4}
% derived from the macros above (spdVsixPack)
\newcommand{\spdVsixPack}{1.4}
% derived from the macros above (spdSFive)
\newcommand{\spdSFive}{6.4}
% derived from the macros above (spdSEight)
\newcommand{\spdSEight}{9.0}
% derived from the macros above (sigRatioHL)
\newcommand{\sigRatioHL}{1.40}
% derived from the macros above (pvfOverTgtHL)
\newcommand{\pvfOverTgtHL}{1.5}
% derived from the macros above (pvfOverTgtRunTwo)
\newcommand{\pvfOverTgtRunTwo}{2.0}
% MANUAL [density]: PI decision 2026-11-11; 3 x 56.45 um, matching_window_audit_2026-11-09.md
\newcommand{\windowMM}{0.17}
% MANUAL [density]: ITk paper Sec. 5.5 hard-scatter criterion
\newcommand{\hsWindowMM}{0.1}
% MANUAL [density]: matching_window_audit_2026-11-09.md (55.05 +- 0.23 um)
\newcommand{\pSixEightPVF}{55.1}
% MANUAL [density]: matching_window_audit_2026-11-09.md (56.45 +- 0.26 um)
\newcommand{\pSixEightAMVF}{56.5}
% MANUAL [density]: note Sec. 3 / v6 config: million training events
\newcommand{\nTrainEvt}{5.17}
% MANUAL [density]: v6 parameter count 3,547,993 (cpu_bench.json), millions
\newcommand{\nParamVsix}{3.55}
% MANUAL [density]: note Table tab:athena-timing-breakdown, CPU total, Xeon 4216
\newcommand{\athCPUms}{2401}
% MANUAL [density]: note Table tab:athena-timing-breakdown, GPU total, T4
\newcommand{\athGPUms}{76}
% MANUAL [density]: note Table tab:athena-timing-breakdown, CPU inference
\newcommand{\athCPUinf}{2189}
% MANUAL [gnn]: note Table tab:athena-timing-breakdown, CPU track association
\newcommand{\athGNNcpu}{182}
% MANUAL [density]: read off note Fig. athena_cpu_gpu_amvf at mu = 200 (600-750 ms)
\newcommand{\athAMVF}{650}
% MANUAL [density]: note Sec. 5.1.2: GPU chain vs AMVF CPU at mu = 200
\newcommand{\athGPUfactor}{9}
% MANUAL [density]: projection 2189/6.4 + 31 + 182 ms (cpu_small_model_2026-11-09.md)
\newcommand{\projSFive}{555}
% MANUAL [density]: projection 2189/9.0 + 31 + 182 ms
\newcommand{\projSEight}{455}
% MANUAL [gnn]: note Sec. 7 GNN parameter count (v4); update for v2
\newcommand{\gnnParams}{39{,}521}
% MANUAL [gnn]: note Table tab:gnn_training (v4); update for v2
\newcommand{\gnnTrainEvt}{180{,}000}
% MANUAL [gnn]: GNN score cut used for every after-GNN number (SOURCES gnn t)
\newcommand{\gnnT}{0.98}


\newcommand{\img}[2]{\includegraphics[width=\linewidth,height=#1\textheight,keepaspectratio]{images/#2}}

\begin{document}
\framelogo{}

\begin{frame}[plain]\titlepage\end{frame}

%----------------------------------------------------------------------
\begin{frame}{Motivation}
\begin{columns}[c]\begin{column}{.55\linewidth}
\textbf{Find the primary vertices and their tracks at the HL-LHC.}
\par\vspace{1em}
\begin{itemize}\setlength\itemsep{.7em}
\item About $200$ interactions per crossing along a few cm of beam line.
\item PV-Finder: a neural network turns tracks into a vertex-density histogram.
\item A graph neural network (GNN) assigns tracks to the vertices.
\item Compared with AMVF on the same events, with the same definitions.
\end{itemize}
\end{column}\begin{column}{.42\linewidth}\centering
\bigvalue{$\langle\mu\rangle=200$}{HL-LHC, ITk}
\par\vspace{1.2em}
\begin{enumerate}\setlength\itemsep{.35em}\small
\item[\textcolor{pennred}{\textbf{I.}}] Method
\item[\textcolor{pennred}{\textbf{II.}}] Vertex finding at PU200
\item[\textcolor{pennred}{\textbf{III.}}] Track association and hard scatter
\item[\textcolor{pennred}{\textbf{IV.}}] Low pile-up and collision data
\item[\textcolor{pennred}{\textbf{V.}}] Timing and deployment
\end{enumerate}
\end{column}\end{columns}
\end{frame}

\partdivider{I}{Method}{Tracks, network, peaks and track association}

%----------------------------------------------------------------------
\begin{frame}{From Tracks to a Vertex Histogram}
\relax
{\centering\resizebox{.95\linewidth}{!}{\begin{tikzpicture}[x=1cm,y=1cm]
  \tikzset{
    term/.style={pvbox, draw=black!35, fill=black!4, minimum width=2.2cm, minimum height=1.35cm, font=\small},
    net/.style ={pvbox, draw=mydarkblue!55, fill=mydarkblue!10, minimum width=2.2cm, minimum height=1.35cm, font=\small},
    map/.style ={pvbox, draw=myorange!70!black, fill=myorange!14, minimum width=2.2cm, minimum height=1.35cm, font=\small},
    sub/.style ={font=\scriptsize, text=black!50, inner sep=1.5pt}}
  \fill[pennred!4, rounded corners=12pt] (1.45,-1.20) rectangle (9.65,1.20);
  \draw[pennred!35, dashed, line width=0.9pt, rounded corners=12pt] (1.45,-1.20) rectangle (9.65,1.20);
  \node[fill=white, inner sep=2.5pt, font=\scriptsize, text=pennred] at (5.55,1.20) {one network, trained end to end};
  \node[term] (tracks) at (0,0)    {Track\\features};
  \node[net]  (mlp)    at (2.85,0) {Masked\\MLP};
  \node[map]  (lat)    at (5.55,0) {Latent\\map};
  \node[net]  (unet)   at (8.25,0) {1-D\\U-Net};
  \node[map]  (out)    at (11.10,0){Vertex\\histogram};
  \node[term] (pv)     at (14.60,0){Primary\\vertices};
  \node[sub, below=2pt] at (tracks.south) {$(7,\,N_\text{trk})$};
  \node[sub, below=2pt] at (out.south)    {$1{,}000$ bins / interval};
  \node[sub, below=2pt] at (pv.south)     {$z$ positions};
  \draw[pvflow] (tracks) -- (mlp); \draw[pvflow] (mlp) -- (lat);
  \draw[pvflow] (lat) -- (unet);   \draw[pvflow] (unet) -- (out);
  \draw[pvflow, draw=pennred!80] (out) -- (pv);
  \node[font=\scriptsize, text=pennred, align=center, anchor=south] at (12.85,0.14) {peak\\finding};
\end{tikzpicture}}\par}
\par\vspace{1em}
\begin{columns}[t]
\begin{column}{.31\linewidth}\textbf{Inputs}\\[.5em]\small
$d_0,z_0,\sigma_{d_0},\sigma_{z_0},\mathrm{cov}(d_0,z_0)$ and the interval bounds.
\end{column}
\begin{column}{.31\linewidth}\textbf{Network}\\[.5em]\small
Per-track MLP, summed over tracks, then a U-Net. Independent of track order.
\end{column}
\begin{column}{.32\linewidth}\textbf{Output}\\[.5em]\small
Twelve $40$ mm intervals: $12{,}000$ bins of $0.04$ mm over $|z|<240$ mm.
\end{column}\end{columns}
\end{frame}

%----------------------------------------------------------------------
\begin{frame}{Peak Finding and the Working Point}
\begin{columns}[c]\begin{column}{.5\linewidth}\centering
\img{.55}{method_example_e2e.png}
\end{column}\begin{column}{.47\linewidth}
\small
\textbf{One configuration for all HL-LHC results}
\par\vspace{.6em}\centering
\begin{tabular}{lr}\toprule
Amplitude threshold & $0.01$\\
Minimum width & $3$ bins\\
Minimum integral & $0.30$\\
Minimum height & $0.03$\\
Position & local centroid\\
GNN score cut $t$ & $\gnnT$\\
Matching window & $\windowMM$ mm\\
Hard-scatter window & $\hsWindowMM$ mm\\\bottomrule
\end{tabular}
\par\vspace{.6em}\raggedright
Peaks split at local minima; centroid over $\pm3$ bins of the maximum.
\end{column}\end{columns}
\end{frame}

%----------------------------------------------------------------------
\begin{frame}{Track-to-Vertex Association with a GNN}
\begin{columns}[c]\begin{column}{.52\linewidth}
\small
\textcolor{mydarkblue}{\textbf{Track nodes:}} $d_0,z_0$, errors, $\theta,\phi,p_T$.\\[.6em]
\textcolor{myorange}{\textbf{Vertex nodes:}} PV-Finder position and peak height.\\[.6em]
\textbf{Edges:} each track to its $k=20$ nearest vertices, with
\[\frac{z_0-z_{\rm PV}}{\sqrt{\sigma_{z_0}^2+\sigma_{\rm PV}^2}},\quad \frac{d_0}{\sigma_{d_0}},\quad |z_0-z_{\rm PV}|.\]
\textbf{Output:} a score per edge. A track goes to its best vertex if the score passes $t=\gnnT$.\\[.6em]
Graph attention network, $\gnnParams$ parameters.
\end{column}\begin{column}{.45\linewidth}\centering\resizebox{.85\linewidth}{!}{\begin{tikzpicture}[x=0.85cm, y=0.85cm]
  \foreach \y/\n in {3.6/1, 2.4/2, 1.2/3, 0.0/4} {
    \filldraw[fill=mydarkblue!14, draw=mydarkblue!70, thick] (0,\y) circle (8pt) node[font=\scriptsize] {\(t_\n\)};}
  \foreach \y/\n in {3.0/1, 1.8/2, 0.6/3} {
    \filldraw[fill=myorange!20, draw=myorange!80!black, thick] (4,\y) circle (8pt) node[font=\scriptsize] {\(v_\n\)};}
  \foreach \s/\t in {3.6/3.0, 3.6/1.8, 2.4/3.0, 2.4/1.8, 1.2/1.8, 1.2/0.6, 0.0/1.8, 0.0/0.6} {
    \draw[black!40, thick] (0.3,\s) -- (3.7,\t);}
  \draw[pennred, very thick] (0.3,2.4) -- (3.7,3.0);
  \node[font=\tiny, text=pennred] at (2.0,3.35) {edge score \(\in[0,1]\)};
  \node[font=\tiny, text=black!55] at (0.0,-0.75) {tracks};
  \node[font=\tiny, text=black!55] at (4.0,-0.75) {PV-Finder vertices};
\end{tikzpicture}}\end{column}\end{columns}
\end{frame}

%----------------------------------------------------------------------
\begin{frame}{HL-LHC Samples}
\begin{columns}[c]\begin{column}{.5\linewidth}
\textbf{$\sqrt{s}=14$ TeV $t\bar t$, ITk, $|\eta|<4$}
\par\vspace{.8em}
\begin{itemize}\setlength\itemsep{.6em}\small
\item Density network: $\nTrainEvt$M training events near $\mu=200$.
\item GNN: $\gnnTrainEvt$ training events.
\item Evaluation on held-out files, flat in $\mu$ from $0$ to $200$.
\item PU200 results: $\nEvtPU$ events at $185\leq\mu\leq200$, $\langle\mu\rangle=\meanMu$.
\item $\truthPerEvt$ truth vertices per event with $\geq2$ tracks.
\end{itemize}
\end{column}\begin{column}{.47\linewidth}\centering
\img{.55}{mu_distribution.pdf}
\end{column}\end{columns}
\context{Tracks: $p_T\geq0.9$ GeV, $|d_0|\leq1$ mm, $\sigma(d_0)\leq0.35$ mm, $|z_0|\leq200$ mm.}
\end{frame}

%----------------------------------------------------------------------
\begin{frame}{Matching Window: $\windowMM$ mm}
\begin{columns}[c]\begin{column}{.55\linewidth}\centering
\img{.6}{window_residuals.png}
\end{column}\begin{column}{.42\linewidth}
\small\textbf{Single-vertex position residual}
\par\vspace{.6em}
$68\%$ of matched vertices lie within:
\par\vspace{.5em}\centering
\begin{tabular}{lr}\toprule
PV-Finder & $\pSixEightPVF~\mu$m\\
AMVF & $\pSixEightAMVF~\mu$m\\\bottomrule
\end{tabular}
\par\vspace{.7em}\raggedright
Window $=3\times\pSixEightAMVF~\mu\mathrm{m}\simeq\windowMM$ mm, the same for both algorithms.
\par\vspace{.5em}
Hard scatter: $\hsWindowMM$ mm, as in the ITk performance studies.
\end{column}\end{columns}
\end{frame}

%----------------------------------------------------------------------
\begin{frame}{Vertex Categories}
\context{Match reconstructed to truth vertices by $z$, closest first, one to one.}
{\centering\resizebox{.95\linewidth}{!}{\begin{tikzpicture}[x=1cm, y=1cm]
  \tikzset{
    truthv/.style={circle, inner sep=0pt, minimum size=9.2pt, fill=black!78, draw=black!88, line width=0.5pt},
    recov/.style ={rectangle, rounded corners=1.5pt, inner sep=0pt, minimum width=8pt, minimum height=10pt, line width=0.6pt},
    band/.style  ={dashed, line width=0.5pt},
    note/.style  ={font=\scriptsize, text=black!55, align=center, anchor=north},
    head/.style  ={font=\small\bfseries}}
  \foreach \cx/\lab/\col in {2.05/Clean/{mygreen!60!black}, 6.05/Merged/{myorange!85!black},
                             10.05/Split/{myblue}, 14.05/Fake/{myred}} {
    \node[pvcard, minimum width=3.5cm, minimum height=3.6cm] at (\cx,1.40) {};
    \node[head, text=\col] at (\cx,2.88) {\lab};
    \draw[\col, line width=1.1pt] (\cx-0.45,2.66) -- (\cx+0.45,2.66);
    \draw[pvguide] (\cx-1.55,2.15) -- (\cx+1.55,2.15);
    \draw[pvguide] (\cx-1.55,0.85) -- (\cx+1.55,0.85);}
  \node[font=\scriptsize, text=black!45, anchor=east] at (0.18,2.15) {truth};
  \node[font=\scriptsize, text=black!45, anchor=east] at (0.18,0.85) {reco};
  \fill[mygreen!10] (1.42,1.85) rectangle (2.82,2.45); \draw[band, draw=mygreen!35] (1.42,1.85) rectangle (2.82,2.45);
  \node[truthv] at (2.00,2.15) {}; \node[recov, fill=mygreen!45, draw=mygreen!65!black] at (2.12,0.85) {};
  \draw[pvpoint, draw=mygreen!60!black] (2.12,1.03) -- (2.02,1.93);
  \node[note] at (2.05,0.35) {one reco,\\one truth};
  \fill[myorange!10] (5.15,1.85) rectangle (6.55,2.45); \draw[band, draw=myorange!40] (5.15,1.85) rectangle (6.55,2.45);
  \node[truthv] at (5.72,2.15) {}; \node[truthv] at (6.42,2.15) {};
  \node[recov, fill=myorange!55, draw=myorange!85!black] at (5.85,0.85) {};
  \draw[pvpoint, draw=myorange!85!black] (5.85,1.03) -- (5.75,1.93);
  \draw[pvpoint, draw=myorange!85!black, dash pattern=on 2.4pt off 1.8pt] (5.97,1.03) -- (6.35,1.93);
  \node[note] at (6.05,0.35) {the second truth\\is credited too};
  \fill[myblue!8] (9.92,1.85) rectangle (11.32,2.45); \draw[band, draw=myblue!30] (9.92,1.85) rectangle (11.32,2.45);
  \node[truthv] at (10.00,2.15) {};
  \node[recov, fill=mygreen!45, draw=mygreen!65!black] at (9.80,0.85) {};
  \node[recov, fill=myblue!45, draw=myblue!80!black] at (10.62,0.85) {};
  \draw[pvpoint, draw=mygreen!60!black] (9.80,1.03) -- (9.94,1.93);
  \draw[myblue!70, dotted, line width=1.0pt] (10.62,1.03) -- (10.16,1.85);
  \node[font=\small, text=myblue!80!black] at (10.52,1.53) {$\times$};
  \node[note] at (10.05,0.35) {a closer reco\\already took it};
  \fill[myred!8] (13.75,1.85) rectangle (15.15,2.45); \draw[band, draw=myred!35] (13.75,1.85) rectangle (15.15,2.45);
  \node[truthv] at (13.05,2.15) {};
  \node[recov, fill=myred!50, draw=myred!80!black] at (14.45,0.85) {};
  \draw[myred!65, dotted, line width=1.0pt] (14.45,1.03) -- (14.45,1.92);
  \node[font=\small, text=myred] at (14.45,2.15) {$\times$};
  \node[note] at (14.05,0.35) {no truth inside\\the window};
\end{tikzpicture}}\par}
\par\vspace{.6em}\small
\textbf{Efficiency:} truth vertices with a reconstructed vertex within the window (clean or merged) / all truth vertices with $\geq2$ tracks.
\end{frame}

\partdivider{II}{Vertex finding at PU200}{Before track association, same events as AMVF}

%----------------------------------------------------------------------
\begin{frame}{PU200: PV-Finder against AMVF}
\begin{columns}[c]\begin{column}{.55\linewidth}\centering
\img{.62}{categories_pu200.png}
\end{column}\begin{column}{.42\linewidth}
\centering\small
\begin{tabular}{lrr}\toprule
 & PV-Finder & AMVF\\\midrule
Efficiency [\%] & $\effPVF$ & $\effAMVF$\\
Fake / event & $\fakePVF$ & $\fakeAMVF$\\
Split / event & $\splitPVF$ & $\splitAMVF$\\\bottomrule
\end{tabular}
\par\vspace{1em}
\bigvalue{$+\effDiff\pm\effDiffErr$}{points of efficiency over AMVF}
\par\vspace{.8em}\raggedright\footnotesize
Hard-scatter window ($\hsWindowMM$ mm): $\effHSwPVF\%$ vs $\effHSwAMVF\%$.
\end{column}\end{columns}
\end{frame}

%----------------------------------------------------------------------
\begin{frame}{Vertex Yield across Pile-up}
\begin{columns}[c]\begin{column}{.62\linewidth}\centering
\img{.68}{categories_vs_mu.png}
\end{column}\begin{column}{.35\linewidth}
\small
\textbf{Reconstructed vertices per event, by category}
\par\vspace{.8em}
Clean vertices dominate at every $\mu$ up to $200$.
\end{column}\end{columns}
\end{frame}

%----------------------------------------------------------------------
\begin{frame}{Vertex--Vertex Resolution}
\begin{columns}[c]\begin{column}{.5\linewidth}\centering
\img{.55}{deltaz_pu200.png}\\[.2em]\footnotesize $\Delta z$ between reconstructed vertices, PU200
\end{column}\begin{column}{.47\linewidth}\centering
\img{.55}{sigma_vs_mu.png}\\[.2em]\footnotesize $\svtx$ versus $\mu$
\end{column}\end{columns}
\par\vspace{.4em}
\takeaway{PU200: $\svtx=\sigPVF\pm\sigPVFErr$ mm for PV-Finder, $\sigAMVF\pm\sigAMVFErr$ mm for AMVF.}
\end{frame}

%----------------------------------------------------------------------
\begin{frame}{Hard-Scatter Vertex Reconstruction}
\begin{columns}[c]\begin{column}{.6\linewidth}\centering
\img{.66}{hs_reco_eff_vs_mu.png}
\end{column}\begin{column}{.37\linewidth}
\small A reconstructed vertex within $\hsWindowMM$ mm of the hard-scatter vertex.
\par\vspace{1.2em}\centering
\bigvalue{$\hsRecoPVF\%$}{PV-Finder at PU200}
\par\vspace{.8em}
{\Large $\hsRecoAMVF\%$}\\[.3em]\small AMVF
\end{column}\end{columns}
\end{frame}

\partdivider{III}{Track association and hard scatter}{After the GNN: categories, hard scatter and position}

%----------------------------------------------------------------------
\begin{frame}{Full Chain at PU200}
\begin{columns}[c]\begin{column}{.6\linewidth}\centering
\img{.66}{clean_yield_pu200.png}
\end{column}\begin{column}{.37\linewidth}
\small\textbf{Clean vertices / truth vertices}
\par\vspace{.6em}
Clean: $\geq70\%$ of the vertex tracks from one interaction.
\par\vspace{.8em}\centering
\begin{tabular}{lr}\toprule
PV-Finder + GNN & $\cleanTruthGNN\%$\\
Nearest $z_0$ & $\cleanTruthZnear\%$\\
$z_0$ within $3\sigma$ & $\cleanTruthZsig\%$\\
AMVF & $\cleanTruthAMVF\%$\\\bottomrule
\end{tabular}
\end{column}\end{columns}
\takeaway{The GNN more than doubles the clean-vertex yield of AMVF and of $z_0$-based assignment.}
\end{frame}

%----------------------------------------------------------------------
\begin{frame}{Full Chain across Pile-up}
\begin{columns}[c]\begin{column}{.65\linewidth}\centering
\img{.7}{clean_yield_vs_mu.png}
\end{column}\begin{column}{.32\linewidth}
\small Clean vertices / truth vertices versus $\mu$ for each stage of the chain and for AMVF.
\par\vspace{.8em}
Lower panel: ratio to AMVF.
\end{column}\end{columns}
\end{frame}

%----------------------------------------------------------------------
\begin{frame}{Hard-Scatter Vertex after Track Association}
\begin{columns}[c]\begin{column}{.47\linewidth}
\centering\small
\begin{tabular}{lrr}\toprule
PU200 & GNN & AMVF\\\midrule
Clean HS vertex & $\cleanHSGNN\%$ & $\cleanHSAMVF\%$\\
HS tracks on HS vertex & $\hsTrkGNN\%$ & $\hsTrkAMVF\%$\\
HS vertex purity & $\hsPurGNN\%$ & $\hsPurAMVF\%$\\
HS selected ($\sum p_T^2$) & $\hsSelGNN\%$ & $\hsSelAMVF\%$\\\bottomrule
\end{tabular}
\par\vspace{.8em}\raggedright\footnotesize
ATLAS definitions; tracks counted, not fit-weighted. HS vertex: most HS tracks. Selected: largest $\sum p_T^2$ and within $\hsWindowMM$ mm. AMVF: its stored primary vertex.
\end{column}\begin{column}{.5\linewidth}\centering
\img{.6}{hs_selection_eff_vs_mu.png}
\end{column}\end{columns}
\end{frame}

%----------------------------------------------------------------------
\begin{frame}{Hard-Scatter Longitudinal Resolution}
\begin{columns}[c]\begin{column}{.55\linewidth}\centering
\img{.66}{hs_resolution_vs_ntrk.png}
\end{column}\begin{column}{.42\linewidth}
\small
\textbf{Vertex fit on the GNN tracks}
\par\vspace{.5em}
Adaptive 3D fit with the beam-spot constraint, seeded at the PV-Finder peak.
\par\vspace{.5em}
ITk definition: iterative RMS of $z_{\rm reco}-z_{\rm HS}$, $\hsResNcommon$ events with the HS found by both.
\par\vspace{.8em}\centering
\bigvalue{$\hsResGNN\pm\hsResGNNErr~\mu$m}{PV-Finder + GNN}
\par\vspace{.5em}
{\Large $\hsResAMVF\pm\hsResAMVFErr~\mu$m}\\[.3em]\small AMVF
\end{column}\end{columns}
\end{frame}

\partdivider{IV}{Low pile-up and collision data}{Run 2 simulation, Run 2 and Run 3 data}

%----------------------------------------------------------------------
\begin{frame}{Low Pile-up: Simulation and Collision Data}
\begin{columns}[c]\begin{column}{.49\linewidth}\centering
\img{.5}{datamc_run2_resolution.png}\\[.2em]\footnotesize Run 2 data and simulation
\end{column}\begin{column}{.49\linewidth}\centering
\img{.5}{datamc_run3_resolution.png}\\[.2em]\footnotesize Run 3 data and simulation
\end{column}\end{columns}
\par\vspace{.3em}\centering\footnotesize
\begin{tabular}{lccc}\toprule
 & Simulation & Run 2 data & Run 3 data\\\midrule
PV-Finder $\svtx$ [mm] & $\sigMC$ & $\sigRunTwo$ & $\sigRunThree$\\\bottomrule
\end{tabular}
\par\vspace{.2em}
Simulation: efficiency $\effMC\%$, $\noRefMCtruth$ fake vertices per event; AMVF $\svtx=\sigMCAMVF$ mm.
\end{frame}

\partdivider{V}{Timing and deployment}{Athena integration, CPU and GPU}

%----------------------------------------------------------------------
\begin{frame}{Athena: Timing on CPU and GPU}
\begin{columns}[c]\begin{column}{.58\linewidth}\centering
\img{.62}{athena_cpu_gpu_amvf.png}
\end{column}\begin{column}{.39\linewidth}
\small
\textbf{Same events, tracks to vertices with track links}
\par\vspace{.6em}
Intel Xeon Silver 4216, NVIDIA T4.
\par\vspace{.6em}\centering
\begin{tabular}{lr}\toprule
At $\mu=200$ & ms/event\\\midrule
PV-Finder + GNN, GPU & $\athGPUms$\\
PV-Finder + GNN, CPU & $\athCPUms$\\
AMVF, CPU & $\sim\athAMVF$\\\bottomrule
\end{tabular}
\par\vspace{.6em}
\textcolor{pennred}{GPU chain $\sim\athGPUfactor\times$ faster than AMVF.}
\end{column}\end{columns}
\context{Implemented in Athena by Trey Warner. Tools and validation in the backup.}
\end{frame}

%----------------------------------------------------------------------
\begin{frame}{A Smaller Network for CPU}
\begin{columns}[t]\begin{column}{.52\linewidth}
\small\textbf{Same data and loss, narrower U-Net}
\par\vspace{.5em}\centering
\begin{tabular}{lrrr}\toprule
 & v6 & s5 & s8\\\midrule
U-Net channels & 280 & 112 & 80\\
CPU time [ms] & $\cpuVsixPad$ & $\cpuSFivePack$ & $\cpuSEightPack$\\
Speed-up & $1$ & $\spdSFive\times$ & $\spdSEight\times$\\
$\Delta$efficiency [pt] & & $\sFiveDEffA$ & $\sEightDEffA$\\
$\Delta$fake / event & & $\sFiveDFake$ & $\sEightDFake$\\\bottomrule
\end{tabular}
\par\vspace{.5em}\raggedright\footnotesize
ONNX Runtime, one thread, one PU200 event. Small models use packed inference: tracks only, no padding, identical output.
\end{column}\begin{column}{.45\linewidth}
\small\textbf{Projected Athena time at $\mu=200$, CPU}
\par\vspace{.5em}\centering
\begin{tabular}{lr}\toprule
Chain & ms/event\\\midrule
v6 + GNN & $\athCPUms$\\
s5 + GNN & $\sim\projSFive$\\
s8 + GNN & $\sim\projSEight$\\
AMVF & $\sim\athAMVF$\\\bottomrule
\end{tabular}
\par\vspace{.6em}\raggedright\footnotesize
Packed inference alone: $\spdVsixPack\times$ for v6 at no cost.
\end{column}\end{columns}
\takeaway{The small network brings the CPU chain to the speed of AMVF at $\mu=200$.}
\end{frame}

%----------------------------------------------------------------------
\begin{frame}{Summary}
\begin{itemize}\setlength\itemsep{.75em}
\item At PU200, PV-Finder finds $\effPVF\%$ of truth vertices, $+\effDiff$ points over AMVF, with fewer fakes.
\item Vertex--vertex resolution $\sigPVF$ mm against $\sigAMVF$ mm for AMVF.
\item With GNN track association: $\cleanTruthGNN\%$ clean vertices per truth vertex, against $\cleanTruthAMVF\%$ for AMVF.
\item Hard scatter: selected in $\hsSelGNN\%$ of events ($\hsSelAMVF\%$ AMVF); position resolution $\hsResGNN~\mu$m ($\hsResAMVF~\mu$m AMVF).
\item Validated on Run 2 and Run 3 collision data.
\item Running in Athena on CPU and GPU.
\end{itemize}
\par\vspace{.6em}
\textbf{Next:} freeze the HL-LHC density model and GNN, and complete the paper draft.
\end{frame}

%----------------------------------------------------------------------
\backupbegin

\begin{frame}{Backup: Efficiency Difference versus Matching Window}
\begin{columns}[c]\begin{column}{.6\linewidth}\centering
\img{.68}{window_scan_eff_diff.png}
\end{column}\begin{column}{.37\linewidth}
\small PV-Finder minus AMVF efficiency, same events, for windows from $0.1$ to $1$ mm.
\par\vspace{.8em}
PV-Finder is ahead at every window.
\end{column}\end{columns}
\end{frame}

\begin{frame}{Backup: Run 2 and HL-LHC Resolution}
\begin{columns}[c]\begin{column}{.52\linewidth}\centering
\img{.6}{sigma_summary.png}
\end{column}\begin{column}{.45\linewidth}
\small
\textbf{AMVF / PV-Finder $\svtx$:} $\ratioRunTwo$ in Run 2 simulation, $\sigRatioHL$ at PU200.
\par\vspace{.6em}
Track-based vertex precision: $\idealPURunTwo~\mu$m (Run 2), $\idealPUHL~\mu$m (HL-LHC).
\par\vspace{.6em}
\textbf{Training target through the same peak finding:}
\par\vspace{.3em}\centering\footnotesize
\begin{tabular}{lrr}\toprule
$\svtx$ [mm] & Run 2 & PU200\\\midrule
Target & $\tgtSigRunTwo$ & $\tgtSigHL$\\
PV-Finder & $\sigMCPVFr$ & $\sigPVF$\\
PV-Finder / target & $\pvfOverTgtRunTwo$ & $\pvfOverTgtHL$\\\bottomrule
\end{tabular}
\par\vspace{.4em}\raggedright
A $0.02$ mm grid moves the target from $\tgtSigHLwone$ to $\tgtSigHLfine$ mm.
\end{column}\end{columns}
\end{frame}

\begin{frame}{Backup: AMVF Hard-Scatter Selection}
\begin{columns}[c]\begin{column}{.55\linewidth}\centering
\img{.62}{amvf_hs_selection_vs_mu.png}
\end{column}\begin{column}{.42\linewidth}
\small
\textbf{AMVF stored primary vertex:} $\hsSelAMVF\%$ at PU200, as in the ITk performance paper.
\par\vspace{.6em}
Re-ranking the stored track lists by $\sum p_T^2$ gives $\hsSelAMVFrerank\%$: a few very high-$p_T$ tracks change the leading vertex.
\end{column}\end{columns}
\end{frame}

\begin{frame}{Backup: Hard-Scatter Tracks and Clean HS versus Pile-up}
\begin{columns}[c]\begin{column}{.49\linewidth}\centering
\img{.62}{hs_clean_fraction_vs_mu.png}
\end{column}\begin{column}{.49\linewidth}\centering
\img{.62}{hs_track_fraction_vs_mu.png}
\end{column}\end{columns}
\end{frame}

\begin{frame}{Backup: Low Pile-up Vertex Yield and the Reference}
\begin{columns}[c]\begin{column}{.52\linewidth}\centering
\img{.62}{datamc_run2_yield.png}
\end{column}\begin{column}{.45\linewidth}
\small
\textbf{PV-Finder peaks without a reference vertex, per event}
\par\vspace{.5em}\centering
\begin{tabular}{lr}\toprule
Simulation, truth reference & $\noRefMCtruth$\\
Simulation, AMVF reference & $\noRefMCamvf$\\
Run 2 data, AMVF reference & $\noRefRunTwo$\\
Run 3 data, AMVF reference & $\noRefRunThree$\\\bottomrule
\end{tabular}
\par\vspace{.6em}\raggedright
In data the reference is AMVF; simulation with the same reference reproduces the data.
\end{column}\end{columns}
\end{frame}

\begin{frame}{Backup: Athena Integration}
\relax
{\centering\resizebox{.94\linewidth}{!}{%
\begin{tikzpicture}[x=1cm,y=1cm,font=\small]
\foreach \x/\id/\txt in {0/trk/{xAOD tracks},3.4/sel/{Track selection\\+ tensor building},6.8/net/{ONNX inference\\+ peak finding},10.2/vtx/{xAOD vertices}}{
\node[pvbox,draw=mydarkblue!60,fill=mydarkblue!7,minimum width=2.9cm,minimum height=1.1cm] (\id) at (\x,0) {\txt};}
\draw[pvflow] (trk)--(sel);\draw[pvflow] (sel)--(net);\draw[pvflow] (net)--(vtx);
\node[pvbox,draw=pennred!70,fill=pennred!7,text width=5.0cm,minimum height=1.0cm] (gnn) at (6.8,-2.0) {GNN: build graph $\to$ score edges\\$\to$ select track links};
\draw[pvflow] (trk.south) -- (0,-.95) -| ([xshift=-12mm]gnn.north);
\draw[pvflow] (vtx.south) -- (10.2,-2.0) -- (gnn.east);
\node[pvbox,draw=mygreen!60!black,fill=mygreen!7,text width=3.1cm,minimum height=1.0cm] (links) at (0,-2.0) {Vertices with\\linked tracks};
\draw[pvflow,draw=mygreen!60!black] (gnn.west)--(links.east);
\end{tikzpicture}}\par}
\par\vspace{.8em}\small
\texttt{InDetNNPriVxFinderTool} and \texttt{InDetGNNTrackVertexAssociationTool}, with the same \texttt{IVertexFinder} interface as AMVF. ONNX Runtime on CPU or CUDA, set by configuration. Athena output agrees with the Python implementation to floating-point precision.
\par\vspace{.5em}{\scriptsize\color{black!55}\treysource.}
\end{frame}

\backupend
\end{document}
")
%
% \readingonly{...} keeps its argument only in the reading build; \talkonly{...}
% only in the talk build.  Every number lives in exactly one place, so the two
% versions cannot drift apart.
\providecommand{\talkbuild}{0}
\newcommand{\readingonly}[1]{\ifnum\talkbuild=0\relax#1\fi}
\newcommand{\talkonly}[1]{\ifnum\talkbuild=1\relax#1\fi}
% \bnum{reading}{talk} expands to a bare number, so it can be used inside a
% length: \includegraphics[height=\bnum{0.5}{0.6}\textheight]{...}.  The talk
% build drops prose, which frees vertical room for the figure that is left.
\newcommand{\bnum}[2]{\ifnum\talkbuild=0 #1\else #2\fi}
% ---------------------------------------------------------------------------

% Part-divider helper: \partdivider{numeral}{title}{one-line subtitle}
\newcommand{\partdivider}[3]{%
    \begin{frame}[plain]
        \centering
        \vfill
        \begin{tikzpicture}
            \fill[pennred] (0,0) circle (0.72cm);
            \node[text=white, font=\Large\bfseries] at (0,0.02) {#1};
        \end{tikzpicture}\\[0.85em]
        {\LARGE\textbf{\textcolor{pennred}{#2}}}\\[0.75em]
        {\color{pennred!28}\rule{0.30\linewidth}{1.1pt}}\\[0.75em]
        {\small\textcolor{black!50}{#3}}
        \vfill
    \end{frame}%
}

%%%%%%%%%%%%%%%%%%%%%%%%%%%%%%%%%%%%%%

\usetheme{Philadelphia}
\hypersetup{colorlinks,allcolors=pennred}

\title[PV-Finder]{PV-Finder: primary vertex reconstruction and track association with neural networks}
\subtitle{ANA-IDTR-2026-01, Editorial Board}

\author[Mattia Mauro]{Qi-Bin Lei, Rocky B. Garg, Lauren Tompkins,\\\textcolor{pennred}{Mattia Mauro}, Trey Warner}

\institute[]{Stanford University}

\date{November 13, 2026}

\titlegraphic{
    \hspace{5pt}
    \href{https://stanford.edu/group/stanford_atlas/}{\includegraphics[height=1cm]{logo-Stanford/Block_S_2_color.png}}
    \href{https://atlas.cern/}{\includegraphics[trim={0 11.5cm 0 11.5cm},clip,height=1cm]{logo-ATLAS/ATLAS-Logo-Ref-Pantone.pdf}}
    \href{https://iris-hep.org/}{\includegraphics[height=1cm]{logo-iris-hep/Iris-hep-4-no-long-name.png}}
}


\newcommand{\readingnote}[1]{\readingonly{\par\vspace{.35em}%
  {\footnotesize\color{black!65}#1\par}}}
\newcommand{\takeaway}[1]{\par\vspace{.65em}{\centering
  \setlength{\fboxsep}{7pt}\colorbox{pennred!5}{\parbox{.94\linewidth}{%
  \centering\small\textcolor{pennred}{#1}}}\par}}
\newcommand{\context}[1]{{\centering\footnotesize\color{black!55}#1\par}\vspace{.65em}}
\newcommand{\bigvalue}[2]{{\LARGE\bfseries\textcolor{pennred}{#1}}\\[.4em]{\small #2}\par}
\newcommand{\pubnote}{\href{https://cds.cern.ch/record/2858348}{ATL-PHYS-PUB-2023-011}}
\newcommand{\notelink}{\href{https://gitlab.cern.ch/atlas-physics-office/IDTR/ANA-IDTR-2026-01/ANA-IDTR-2026-01-INT1}{ANA-IDTR-2026-01-INT1}}
\newcommand{\treysource}{\href{https://gitlab.cern.ch/atlas-physics-office/IDTR/ANA-IDTR-2026-01/ANA-IDTR-2026-01-INT1/-/blob/4c50ed51500c897dd6a1b5bb52b0f1e3c12e7575/sections/04b_timing.tex}{Note Sec. 5.1 (18 Sep 2026)}}

\newcommand{\svtx}{\sigma_{\mathrm{vtx-vtx}}}
\setbeamertemplate{navigation symbols}{}
\hypersetup{pdftitle={PV-Finder, EB presentation, November 13, 2026},
  pdfauthor={Mattia Mauro, Qi-Bin Lei, Rocky B. Garg, Lauren Tompkins, Trey Warner}}

% numbers.tex: GENERATED by scripts/eb_numbers.py. Do not edit by hand.
%
% [density] checkpoint: hllhc_alleta_v6_mse_2ep_phase2_epoch_2 (v6 production)
%     window: 0.1694 mm (to become 0.17 mm)
%     peaks: amplitude 0.01, width >= 3 bins, integral 0.30, height 0.03, centroid
% [gnn] checkpoint: ttva_gnn_hllhc_v4_noaug epoch 107 (to be replaced by full-mu v2)
%     t: t098
% [lowpu] checkpoint: Run 2 reproduction model (reproduction_ep300), integral 0.50
%
% density: outputs/11_09_2026_output/split_inclusive_efficiency/results.json :: windows/r16443@0.1694/pvf/eff
\newcommand{\effPVF}{83.2}
% density: outputs/11_09_2026_output/split_inclusive_efficiency/results.json :: windows/r16443@0.1694/amvf/eff
\newcommand{\effAMVF}{79.8}
% density: outputs/11_09_2026_output/split_inclusive_efficiency/results.json :: windows/r16443@0.1694/diff/eff
\newcommand{\effDiff}{3.36}
% density: outputs/11_09_2026_output/split_inclusive_efficiency/results.json :: windows/r16443@0.1694/diff/eff_err
\newcommand{\effDiffErr}{0.06}
% density: outputs/11_09_2026_output/split_inclusive_efficiency/results.json :: windows/r16443@0.1/pvf/eff
\newcommand{\effHSwPVF}{74.1}
% density: outputs/11_09_2026_output/split_inclusive_efficiency/results.json :: windows/r16443@0.1/amvf/eff
\newcommand{\effHSwAMVF}{71.0}
% density: outputs/11_09_2026_output/split_inclusive_efficiency/results.json :: windows/r16443@0.1/diff/eff
\newcommand{\effHSwDiff}{3.12}
% density: outputs/11_09_2026_output/split_inclusive_efficiency/results.json :: windows/r16443@0.1694/pvf/fake_per_evt
\newcommand{\fakePVF}{18.7}
% density: outputs/11_09_2026_output/split_inclusive_efficiency/results.json :: windows/r16443@0.1694/amvf/fake_per_evt
\newcommand{\fakeAMVF}{19.7}
% density: outputs/11_09_2026_output/split_inclusive_efficiency/results.json :: windows/r16443@0.1694/pvf/split_per_evt
\newcommand{\splitPVF}{0.53}
% density: outputs/11_09_2026_output/split_inclusive_efficiency/results.json :: windows/r16443@0.1694/amvf/split_per_evt
\newcommand{\splitAMVF}{0.18}
% density: outputs/11_09_2026_output/split_inclusive_efficiency/results.json :: windows/r16443@0.1694/pvf/truth_per_evt
\newcommand{\truthPerEvt}{111.3}
% density: outputs/11_09_2026_output/split_inclusive_efficiency/results.json :: windows/r16443@0.1694/pvf/reco_per_evt
\newcommand{\recoPVF}{101.0}
% density: outputs/11_09_2026_output/split_inclusive_efficiency/results.json :: windows/r16443@0.1694/amvf/reco_per_evt
\newcommand{\recoAMVF}{97.9}
% density: outputs/11_09_2026_output/split_inclusive_efficiency/results.json :: windows/r16443@0.1694/n_events
\newcommand{\nEvtPU}{1{,}920}
% density: outputs/11_09_2026_output/split_inclusive_efficiency/results.json :: windows/r16443@0.1694/totals/pvf/n_truth
\newcommand{\nTruthPU}{213{,}738}
% density: outputs/11_09_2026_output/run2_vs_run4_resolution/results.json :: sigma/hl_pu200/pvf_int030_h003_cen/sigma_mm (integral 0.30, same as counting (one threshold set))
\newcommand{\sigPVF}{0.218}
% density: outputs/11_09_2026_output/run2_vs_run4_resolution/results.json :: sigma/hl_pu200/pvf_int030_h003_cen/fit_err_mm
\newcommand{\sigPVFErr}{0.003}
% density: outputs/11_09_2026_output/run2_vs_run4_resolution/results.json :: sigma/hl_pu200/amvf/sigma_mm
\newcommand{\sigAMVF}{0.305}
% density: outputs/11_09_2026_output/run2_vs_run4_resolution/results.json :: sigma/hl_pu200/amvf/fit_err_mm
\newcommand{\sigAMVFErr}{0.006}
% gnn: outputs/11_09_2026_output/hs_chain_study/hs_chain_summary.json :: reco_min_tracks_1/pu200/fractions/a_reco_pvf_peaks/value
\newcommand{\hsRecoPVF}{99.9}
% gnn: outputs/11_09_2026_output/hs_chain_study/hs_chain_summary.json :: reco_min_tracks_1/pu200/fractions/a_reco_amvf/value
\newcommand{\hsRecoAMVF}{99.8}
% gnn: outputs/11_09_2026_output/hs_chain_study/hs_chain_summary.json :: reco_min_tracks_1/pu200/mean_mu
\newcommand{\meanMu}{192.5}
% density: outputs/11_10_2026_output/small_model_eval/compare3_r16443.json :: s5/d_eff_pt
\newcommand{\sFiveDEffA}{-0.33}
% density: outputs/11_10_2026_output/small_model_eval/compare3_r16638.json :: s5/d_eff_pt
\newcommand{\sFiveDEffB}{-0.26}
% density: outputs/11_10_2026_output/small_model_eval/compare3_r16443.json :: s5/d_fp
\newcommand{\sFiveDFake}{-0.9}
% density: outputs/11_10_2026_output/small_model_eval/compare3_r16443.json :: s8/d_eff_pt
\newcommand{\sEightDEffA}{-0.80}
% density: outputs/11_10_2026_output/small_model_eval/compare3_r16638.json :: s8/d_eff_pt
\newcommand{\sEightDEffB}{-0.70}
% density: outputs/11_10_2026_output/small_model_eval/compare3_r16443.json :: s8/d_fp
\newcommand{\sEightDFake}{-3.1}
% density: outputs/11_09_2026_output/cpu_bench_packed/cpu_bench.json :: 0/median_ms
\newcommand{\cpuVsixPad}{1{,}601}
% density: outputs/11_09_2026_output/cpu_bench_packed/cpu_bench.json :: 1/median_ms
\newcommand{\cpuVsixPack}{1{,}159}
% density: outputs/11_09_2026_output/cpu_bench_packed/cpu_bench.json :: 3/median_ms
\newcommand{\cpuSFivePack}{250}
% density: outputs/11_09_2026_output/cpu_bench_packed/cpu_bench.json :: 5/median_ms
\newcommand{\cpuSEightPack}{178}
% gnn: outputs/11_09_2026_output/hs_chain_study/hs_chain_summary.json :: reco_min_tracks_1/pu200/clean_per_truth/gnn_v4_t098/value
\newcommand{\cleanTruthGNN}{68.4}
% gnn: outputs/11_09_2026_output/hs_chain_study/hs_chain_summary.json :: reco_min_tracks_1/pu200/clean_per_truth/amvf/value
\newcommand{\cleanTruthAMVF}{30.9}
% gnn: outputs/11_09_2026_output/hs_chain_study/hs_chain_summary.json :: reco_min_tracks_1/pu200/clean_per_truth/z0_nearest/value
\newcommand{\cleanTruthZnear}{31.9}
% gnn: outputs/11_09_2026_output/hs_chain_study/hs_chain_summary.json :: reco_min_tracks_1/pu200/clean_per_truth/z0_3sigma/value
\newcommand{\cleanTruthZsig}{33.1}
% gnn: outputs/11_09_2026_output/hs_chain_study/hs_chain_summary.json :: reco_min_tracks_1/pu200/fractions/b_clean_hs_gnn_v4_t098/value
\newcommand{\cleanHSGNN}{96.4}
% gnn: outputs/11_09_2026_output/hs_chain_study/hs_chain_summary.json :: reco_min_tracks_1/pu200/fractions/b_clean_hs_amvf/value
\newcommand{\cleanHSAMVF}{92.5}
% gnn: outputs/11_09_2026_output/hs_chain_study/hs_chain_summary.json :: reco_min_tracks_1/pu200/c_hs_track_fraction/gnn_v4_t098/value
\newcommand{\hsTrkGNN}{74.9}
% gnn: outputs/11_09_2026_output/hs_chain_study/hs_chain_summary.json :: reco_min_tracks_1/pu200/c_hs_track_fraction/amvf/value
\newcommand{\hsTrkAMVF}{74.1}
% gnn: outputs/11_09_2026_output/hs_chain_study/hs_chain_summary.json :: reco_min_tracks_1/pu200/c_hs_vertex_purity/gnn_v4_t098/value
\newcommand{\hsPurGNN}{91.3}
% gnn: outputs/11_09_2026_output/hs_chain_study/hs_chain_summary.json :: reco_min_tracks_1/pu200/c_hs_vertex_purity/amvf/value
\newcommand{\hsPurAMVF}{88.4}
% gnn: outputs/11_09_2026_output/hs_chain_study/hs_chain_summary.json :: reco_min_tracks_1/pu200/fractions/d_select_gnn_v4_t098_itk_0p1/value
\newcommand{\hsSelGNN}{96.3}
% gnn: outputs/11_09_2026_output/hs_chain_study/hs_chain_summary.json :: reco_min_tracks_1/pu200/fractions/d_native_amvf_itk_0p1/value (AMVF stored primary, same code as the GNN number)
\newcommand{\hsSelAMVF}{92.9}
% gnn: outputs/11_09_2026_output/hs_selection_reconciliation/reconciliation.json :: pu200_185_200/flags/rerank_plurality/value
\newcommand{\hsSelAMVFrerank}{80.7}
% gnn: outputs/11_11_2026_output/hs_position_estimators/results.json :: estimators/ada3d_t098/rms_um (3D adaptive fit with beam spot on the GNN tracks at the same t)
\newcommand{\hsResGNN}{9.7}
% gnn: outputs/11_11_2026_output/hs_position_estimators/results.json :: estimators/ada3d_t098/boot_err_um
\newcommand{\hsResGNNErr}{0.3}
% gnn: outputs/11_11_2026_output/hs_position_estimators/results.json :: estimators/amvf/rms_um
\newcommand{\hsResAMVF}{10.2}
% gnn: outputs/11_11_2026_output/hs_position_estimators/results.json :: estimators/amvf/boot_err_um
\newcommand{\hsResAMVFErr}{0.3}
% gnn: outputs/11_11_2026_output/hs_position_estimators/results.json :: estimators/centroid/rms_um
\newcommand{\hsResPeak}{12.4}
% gnn: outputs/11_11_2026_output/hs_position_estimators/results.json :: estimators/ada3d_truth/rms_um
\newcommand{\hsResTruth}{9.1}
% gnn: outputs/11_11_2026_output/hs_position_estimators/results.json :: n_common
\newcommand{\hsResNcommon}{1{,}915}
% lowpu: outputs/11_09_2026_output/data_mc_validation/results.json :: sigma/mc/fit
\newcommand{\sigMC}{0.416}
% lowpu: outputs/11_09_2026_output/data_mc_validation/results.json :: sigma/run2/fit
\newcommand{\sigRunTwo}{0.368}
% lowpu: outputs/11_09_2026_output/data_mc_validation/results.json :: sigma/run3/fit
\newcommand{\sigRunThree}{0.348}
% lowpu: outputs/11_09_2026_output/data_mc_validation/results.json :: matching/fitted_sample/mc_truth_ref/recovery
\newcommand{\effMC}{91.6}
% lowpu: outputs/11_09_2026_output/data_mc_validation/results.json :: matching/fitted_sample/mc_truth_ref/no_ref_per_evt
\newcommand{\noRefMCtruth}{0.42}
% lowpu: outputs/11_09_2026_output/data_mc_validation/results.json :: matching/fitted_sample/mc_amvf_ref/no_ref_per_evt
\newcommand{\noRefMCamvf}{1.72}
% lowpu: outputs/11_09_2026_output/data_mc_validation/results.json :: matching/fitted_sample/run2_data_amvf_ref/no_ref_per_evt
\newcommand{\noRefRunTwo}{1.82}
% lowpu: outputs/11_09_2026_output/data_mc_validation/results.json :: matching/fitted_sample/run3_data_amvf_ref/no_ref_per_evt
\newcommand{\noRefRunThree}{1.96}
% lowpu: outputs/11_09_2026_output/run2_vs_run4_resolution/results.json :: sigma/run2/amvf/sigma_mm
\newcommand{\sigMCAMVF}{0.789}
% lowpu: outputs/11_09_2026_output/run2_vs_run4_resolution/results.json :: sigma/run2/pvf_int050_h000_full/sigma_mm
\newcommand{\sigMCPVFr}{0.410}
% lowpu: outputs/11_09_2026_output/run2_vs_run4_resolution/results.json :: sigma/run2/ratio_amvf_over_pvf/ratio
\newcommand{\ratioRunTwo}{1.92}
% lowpu: outputs/11_11_2026_output/target_readout/results.json :: run2/target_run2cfg/sigma/sigma_mm
\newcommand{\tgtSigRunTwo}{0.203}
% lowpu: outputs/11_11_2026_output/target_readout/results.json :: hl_pu200/target_notecfg/sigma/sigma_mm
\newcommand{\tgtSigHL}{0.145}
% lowpu: outputs/11_11_2026_output/target_readout/results.json :: hl_pu200/target_notecfg_w1/sigma/sigma_mm
\newcommand{\tgtSigHLwone}{0.137}
% lowpu: outputs/11_11_2026_output/target_readout/results.json :: hl_pu200/target_notecfg_w1_g002/sigma/sigma_mm
\newcommand{\tgtSigHLfine}{0.124}
% lowpu: outputs/11_11_2026_output/target_readout/results.json :: ideal_precision_um/run2/pu_median
\newcommand{\idealPURunTwo}{75.6}
% lowpu: outputs/11_11_2026_output/target_readout/results.json :: ideal_precision_um/hl/pu_median
\newcommand{\idealPUHL}{44.4}
% derived from the macros above (spdVsixPack)
\newcommand{\spdVsixPack}{1.4}
% derived from the macros above (spdSFive)
\newcommand{\spdSFive}{6.4}
% derived from the macros above (spdSEight)
\newcommand{\spdSEight}{9.0}
% derived from the macros above (sigRatioHL)
\newcommand{\sigRatioHL}{1.40}
% derived from the macros above (pvfOverTgtHL)
\newcommand{\pvfOverTgtHL}{1.5}
% derived from the macros above (pvfOverTgtRunTwo)
\newcommand{\pvfOverTgtRunTwo}{2.0}
% MANUAL [density]: PI decision 2026-11-11; 3 x 56.45 um, matching_window_audit_2026-11-09.md
\newcommand{\windowMM}{0.17}
% MANUAL [density]: ITk paper Sec. 5.5 hard-scatter criterion
\newcommand{\hsWindowMM}{0.1}
% MANUAL [density]: matching_window_audit_2026-11-09.md (55.05 +- 0.23 um)
\newcommand{\pSixEightPVF}{55.1}
% MANUAL [density]: matching_window_audit_2026-11-09.md (56.45 +- 0.26 um)
\newcommand{\pSixEightAMVF}{56.5}
% MANUAL [density]: note Sec. 3 / v6 config: million training events
\newcommand{\nTrainEvt}{5.17}
% MANUAL [density]: v6 parameter count 3,547,993 (cpu_bench.json), millions
\newcommand{\nParamVsix}{3.55}
% MANUAL [density]: note Table tab:athena-timing-breakdown, CPU total, Xeon 4216
\newcommand{\athCPUms}{2401}
% MANUAL [density]: note Table tab:athena-timing-breakdown, GPU total, T4
\newcommand{\athGPUms}{76}
% MANUAL [density]: note Table tab:athena-timing-breakdown, CPU inference
\newcommand{\athCPUinf}{2189}
% MANUAL [gnn]: note Table tab:athena-timing-breakdown, CPU track association
\newcommand{\athGNNcpu}{182}
% MANUAL [density]: read off note Fig. athena_cpu_gpu_amvf at mu = 200 (600-750 ms)
\newcommand{\athAMVF}{650}
% MANUAL [density]: note Sec. 5.1.2: GPU chain vs AMVF CPU at mu = 200
\newcommand{\athGPUfactor}{9}
% MANUAL [density]: projection 2189/6.4 + 31 + 182 ms (cpu_small_model_2026-11-09.md)
\newcommand{\projSFive}{555}
% MANUAL [density]: projection 2189/9.0 + 31 + 182 ms
\newcommand{\projSEight}{455}
% MANUAL [gnn]: note Sec. 7 GNN parameter count (v4); update for v2
\newcommand{\gnnParams}{39{,}521}
% MANUAL [gnn]: note Table tab:gnn_training (v4); update for v2
\newcommand{\gnnTrainEvt}{180{,}000}
% MANUAL [gnn]: GNN score cut used for every after-GNN number (SOURCES gnn t)
\newcommand{\gnnT}{0.98}


\newcommand{\img}[2]{\includegraphics[width=\linewidth,height=#1\textheight,keepaspectratio]{images/#2}}

\begin{document}
\framelogo{}

\begin{frame}[plain]\titlepage\end{frame}

%----------------------------------------------------------------------
\begin{frame}{Motivation}
\begin{columns}[c]\begin{column}{.55\linewidth}
\textbf{Find the primary vertices and their tracks at the HL-LHC.}
\par\vspace{1em}
\begin{itemize}\setlength\itemsep{.7em}
\item About $200$ interactions per crossing along a few cm of beam line.
\item PV-Finder: a neural network turns tracks into a vertex-density histogram.
\item A graph neural network (GNN) assigns tracks to the vertices.
\item Compared with AMVF on the same events, with the same definitions.
\end{itemize}
\end{column}\begin{column}{.42\linewidth}\centering
\bigvalue{$\langle\mu\rangle=200$}{HL-LHC, ITk}
\par\vspace{1.2em}
\begin{enumerate}\setlength\itemsep{.35em}\small
\item[\textcolor{pennred}{\textbf{I.}}] Method
\item[\textcolor{pennred}{\textbf{II.}}] Vertex finding at PU200
\item[\textcolor{pennred}{\textbf{III.}}] Track association and hard scatter
\item[\textcolor{pennred}{\textbf{IV.}}] Low pile-up and collision data
\item[\textcolor{pennred}{\textbf{V.}}] Timing and deployment
\end{enumerate}
\end{column}\end{columns}
\end{frame}

\partdivider{I}{Method}{Tracks, network, peaks and track association}

%----------------------------------------------------------------------
\begin{frame}{From Tracks to a Vertex Histogram}
\relax
{\centering\resizebox{.95\linewidth}{!}{\begin{tikzpicture}[x=1cm,y=1cm]
  \tikzset{
    term/.style={pvbox, draw=black!35, fill=black!4, minimum width=2.2cm, minimum height=1.35cm, font=\small},
    net/.style ={pvbox, draw=mydarkblue!55, fill=mydarkblue!10, minimum width=2.2cm, minimum height=1.35cm, font=\small},
    map/.style ={pvbox, draw=myorange!70!black, fill=myorange!14, minimum width=2.2cm, minimum height=1.35cm, font=\small},
    sub/.style ={font=\scriptsize, text=black!50, inner sep=1.5pt}}
  \fill[pennred!4, rounded corners=12pt] (1.45,-1.20) rectangle (9.65,1.20);
  \draw[pennred!35, dashed, line width=0.9pt, rounded corners=12pt] (1.45,-1.20) rectangle (9.65,1.20);
  \node[fill=white, inner sep=2.5pt, font=\scriptsize, text=pennred] at (5.55,1.20) {one network, trained end to end};
  \node[term] (tracks) at (0,0)    {Track\\features};
  \node[net]  (mlp)    at (2.85,0) {Masked\\MLP};
  \node[map]  (lat)    at (5.55,0) {Latent\\map};
  \node[net]  (unet)   at (8.25,0) {1-D\\U-Net};
  \node[map]  (out)    at (11.10,0){Vertex\\histogram};
  \node[term] (pv)     at (14.60,0){Primary\\vertices};
  \node[sub, below=2pt] at (tracks.south) {$(7,\,N_\text{trk})$};
  \node[sub, below=2pt] at (out.south)    {$1{,}000$ bins / interval};
  \node[sub, below=2pt] at (pv.south)     {$z$ positions};
  \draw[pvflow] (tracks) -- (mlp); \draw[pvflow] (mlp) -- (lat);
  \draw[pvflow] (lat) -- (unet);   \draw[pvflow] (unet) -- (out);
  \draw[pvflow, draw=pennred!80] (out) -- (pv);
  \node[font=\scriptsize, text=pennred, align=center, anchor=south] at (12.85,0.14) {peak\\finding};
\end{tikzpicture}}\par}
\par\vspace{1em}
\begin{columns}[t]
\begin{column}{.31\linewidth}\textbf{Inputs}\\[.5em]\small
$d_0,z_0,\sigma_{d_0},\sigma_{z_0},\mathrm{cov}(d_0,z_0)$ and the interval bounds.
\end{column}
\begin{column}{.31\linewidth}\textbf{Network}\\[.5em]\small
Per-track MLP, summed over tracks, then a U-Net. Independent of track order.
\end{column}
\begin{column}{.32\linewidth}\textbf{Output}\\[.5em]\small
Twelve $40$ mm intervals: $12{,}000$ bins of $0.04$ mm over $|z|<240$ mm.
\end{column}\end{columns}
\end{frame}

%----------------------------------------------------------------------
\begin{frame}{Peak Finding and the Working Point}
\begin{columns}[c]\begin{column}{.5\linewidth}\centering
\img{.55}{method_example_e2e.png}
\end{column}\begin{column}{.47\linewidth}
\small
\textbf{One configuration for all HL-LHC results}
\par\vspace{.6em}\centering
\begin{tabular}{lr}\toprule
Amplitude threshold & $0.01$\\
Minimum width & $3$ bins\\
Minimum integral & $0.30$\\
Minimum height & $0.03$\\
Position & local centroid\\
GNN score cut $t$ & $\gnnT$\\
Matching window & $\windowMM$ mm\\
Hard-scatter window & $\hsWindowMM$ mm\\\bottomrule
\end{tabular}
\par\vspace{.6em}\raggedright
Peaks split at local minima; centroid over $\pm3$ bins of the maximum.
\end{column}\end{columns}
\end{frame}

%----------------------------------------------------------------------
\begin{frame}{Track-to-Vertex Association with a GNN}
\begin{columns}[c]\begin{column}{.52\linewidth}
\small
\textcolor{mydarkblue}{\textbf{Track nodes:}} $d_0,z_0$, errors, $\theta,\phi,p_T$.\\[.6em]
\textcolor{myorange}{\textbf{Vertex nodes:}} PV-Finder position and peak height.\\[.6em]
\textbf{Edges:} each track to its $k=20$ nearest vertices, with
\[\frac{z_0-z_{\rm PV}}{\sqrt{\sigma_{z_0}^2+\sigma_{\rm PV}^2}},\quad \frac{d_0}{\sigma_{d_0}},\quad |z_0-z_{\rm PV}|.\]
\textbf{Output:} a score per edge. A track goes to its best vertex if the score passes $t=\gnnT$.\\[.6em]
Graph attention network, $\gnnParams$ parameters.
\end{column}\begin{column}{.45\linewidth}\centering\resizebox{.85\linewidth}{!}{\begin{tikzpicture}[x=0.85cm, y=0.85cm]
  \foreach \y/\n in {3.6/1, 2.4/2, 1.2/3, 0.0/4} {
    \filldraw[fill=mydarkblue!14, draw=mydarkblue!70, thick] (0,\y) circle (8pt) node[font=\scriptsize] {\(t_\n\)};}
  \foreach \y/\n in {3.0/1, 1.8/2, 0.6/3} {
    \filldraw[fill=myorange!20, draw=myorange!80!black, thick] (4,\y) circle (8pt) node[font=\scriptsize] {\(v_\n\)};}
  \foreach \s/\t in {3.6/3.0, 3.6/1.8, 2.4/3.0, 2.4/1.8, 1.2/1.8, 1.2/0.6, 0.0/1.8, 0.0/0.6} {
    \draw[black!40, thick] (0.3,\s) -- (3.7,\t);}
  \draw[pennred, very thick] (0.3,2.4) -- (3.7,3.0);
  \node[font=\tiny, text=pennred] at (2.0,3.35) {edge score \(\in[0,1]\)};
  \node[font=\tiny, text=black!55] at (0.0,-0.75) {tracks};
  \node[font=\tiny, text=black!55] at (4.0,-0.75) {PV-Finder vertices};
\end{tikzpicture}}\end{column}\end{columns}
\end{frame}

%----------------------------------------------------------------------
\begin{frame}{HL-LHC Samples}
\begin{columns}[c]\begin{column}{.5\linewidth}
\textbf{$\sqrt{s}=14$ TeV $t\bar t$, ITk, $|\eta|<4$}
\par\vspace{.8em}
\begin{itemize}\setlength\itemsep{.6em}\small
\item Density network: $\nTrainEvt$M training events near $\mu=200$.
\item GNN: $\gnnTrainEvt$ training events.
\item Evaluation on held-out files, flat in $\mu$ from $0$ to $200$.
\item PU200 results: $\nEvtPU$ events at $185\leq\mu\leq200$, $\langle\mu\rangle=\meanMu$.
\item $\truthPerEvt$ truth vertices per event with $\geq2$ tracks.
\end{itemize}
\end{column}\begin{column}{.47\linewidth}\centering
\img{.55}{mu_distribution.pdf}
\end{column}\end{columns}
\context{Tracks: $p_T\geq0.9$ GeV, $|d_0|\leq1$ mm, $\sigma(d_0)\leq0.35$ mm, $|z_0|\leq200$ mm.}
\end{frame}

%----------------------------------------------------------------------
\begin{frame}{Matching Window: $\windowMM$ mm}
\begin{columns}[c]\begin{column}{.55\linewidth}\centering
\img{.6}{window_residuals.png}
\end{column}\begin{column}{.42\linewidth}
\small\textbf{Single-vertex position residual}
\par\vspace{.6em}
$68\%$ of matched vertices lie within:
\par\vspace{.5em}\centering
\begin{tabular}{lr}\toprule
PV-Finder & $\pSixEightPVF~\mu$m\\
AMVF & $\pSixEightAMVF~\mu$m\\\bottomrule
\end{tabular}
\par\vspace{.7em}\raggedright
Window $=3\times\pSixEightAMVF~\mu\mathrm{m}\simeq\windowMM$ mm, the same for both algorithms.
\par\vspace{.5em}
Hard scatter: $\hsWindowMM$ mm, as in the ITk performance studies.
\end{column}\end{columns}
\end{frame}

%----------------------------------------------------------------------
\begin{frame}{Vertex Categories}
\context{Match reconstructed to truth vertices by $z$, closest first, one to one.}
{\centering\resizebox{.95\linewidth}{!}{\begin{tikzpicture}[x=1cm, y=1cm]
  \tikzset{
    truthv/.style={circle, inner sep=0pt, minimum size=9.2pt, fill=black!78, draw=black!88, line width=0.5pt},
    recov/.style ={rectangle, rounded corners=1.5pt, inner sep=0pt, minimum width=8pt, minimum height=10pt, line width=0.6pt},
    band/.style  ={dashed, line width=0.5pt},
    note/.style  ={font=\scriptsize, text=black!55, align=center, anchor=north},
    head/.style  ={font=\small\bfseries}}
  \foreach \cx/\lab/\col in {2.05/Clean/{mygreen!60!black}, 6.05/Merged/{myorange!85!black},
                             10.05/Split/{myblue}, 14.05/Fake/{myred}} {
    \node[pvcard, minimum width=3.5cm, minimum height=3.6cm] at (\cx,1.40) {};
    \node[head, text=\col] at (\cx,2.88) {\lab};
    \draw[\col, line width=1.1pt] (\cx-0.45,2.66) -- (\cx+0.45,2.66);
    \draw[pvguide] (\cx-1.55,2.15) -- (\cx+1.55,2.15);
    \draw[pvguide] (\cx-1.55,0.85) -- (\cx+1.55,0.85);}
  \node[font=\scriptsize, text=black!45, anchor=east] at (0.18,2.15) {truth};
  \node[font=\scriptsize, text=black!45, anchor=east] at (0.18,0.85) {reco};
  \fill[mygreen!10] (1.42,1.85) rectangle (2.82,2.45); \draw[band, draw=mygreen!35] (1.42,1.85) rectangle (2.82,2.45);
  \node[truthv] at (2.00,2.15) {}; \node[recov, fill=mygreen!45, draw=mygreen!65!black] at (2.12,0.85) {};
  \draw[pvpoint, draw=mygreen!60!black] (2.12,1.03) -- (2.02,1.93);
  \node[note] at (2.05,0.35) {one reco,\\one truth};
  \fill[myorange!10] (5.15,1.85) rectangle (6.55,2.45); \draw[band, draw=myorange!40] (5.15,1.85) rectangle (6.55,2.45);
  \node[truthv] at (5.72,2.15) {}; \node[truthv] at (6.42,2.15) {};
  \node[recov, fill=myorange!55, draw=myorange!85!black] at (5.85,0.85) {};
  \draw[pvpoint, draw=myorange!85!black] (5.85,1.03) -- (5.75,1.93);
  \draw[pvpoint, draw=myorange!85!black, dash pattern=on 2.4pt off 1.8pt] (5.97,1.03) -- (6.35,1.93);
  \node[note] at (6.05,0.35) {the second truth\\is credited too};
  \fill[myblue!8] (9.92,1.85) rectangle (11.32,2.45); \draw[band, draw=myblue!30] (9.92,1.85) rectangle (11.32,2.45);
  \node[truthv] at (10.00,2.15) {};
  \node[recov, fill=mygreen!45, draw=mygreen!65!black] at (9.80,0.85) {};
  \node[recov, fill=myblue!45, draw=myblue!80!black] at (10.62,0.85) {};
  \draw[pvpoint, draw=mygreen!60!black] (9.80,1.03) -- (9.94,1.93);
  \draw[myblue!70, dotted, line width=1.0pt] (10.62,1.03) -- (10.16,1.85);
  \node[font=\small, text=myblue!80!black] at (10.52,1.53) {$\times$};
  \node[note] at (10.05,0.35) {a closer reco\\already took it};
  \fill[myred!8] (13.75,1.85) rectangle (15.15,2.45); \draw[band, draw=myred!35] (13.75,1.85) rectangle (15.15,2.45);
  \node[truthv] at (13.05,2.15) {};
  \node[recov, fill=myred!50, draw=myred!80!black] at (14.45,0.85) {};
  \draw[myred!65, dotted, line width=1.0pt] (14.45,1.03) -- (14.45,1.92);
  \node[font=\small, text=myred] at (14.45,2.15) {$\times$};
  \node[note] at (14.05,0.35) {no truth inside\\the window};
\end{tikzpicture}}\par}
\par\vspace{.6em}\small
\textbf{Efficiency:} truth vertices with a reconstructed vertex within the window (clean or merged) / all truth vertices with $\geq2$ tracks.
\end{frame}

\partdivider{II}{Vertex finding at PU200}{Before track association, same events as AMVF}

%----------------------------------------------------------------------
\begin{frame}{PU200: PV-Finder against AMVF}
\begin{columns}[c]\begin{column}{.55\linewidth}\centering
\img{.62}{categories_pu200.png}
\end{column}\begin{column}{.42\linewidth}
\centering\small
\begin{tabular}{lrr}\toprule
 & PV-Finder & AMVF\\\midrule
Efficiency [\%] & $\effPVF$ & $\effAMVF$\\
Fake / event & $\fakePVF$ & $\fakeAMVF$\\
Split / event & $\splitPVF$ & $\splitAMVF$\\\bottomrule
\end{tabular}
\par\vspace{1em}
\bigvalue{$+\effDiff\pm\effDiffErr$}{points of efficiency over AMVF}
\par\vspace{.8em}\raggedright\footnotesize
Hard-scatter window ($\hsWindowMM$ mm): $\effHSwPVF\%$ vs $\effHSwAMVF\%$.
\end{column}\end{columns}
\end{frame}

%----------------------------------------------------------------------
\begin{frame}{Vertex Yield across Pile-up}
\begin{columns}[c]\begin{column}{.62\linewidth}\centering
\img{.68}{categories_vs_mu.png}
\end{column}\begin{column}{.35\linewidth}
\small
\textbf{Reconstructed vertices per event, by category}
\par\vspace{.8em}
Clean vertices dominate at every $\mu$ up to $200$.
\end{column}\end{columns}
\end{frame}

%----------------------------------------------------------------------
\begin{frame}{Vertex--Vertex Resolution}
\begin{columns}[c]\begin{column}{.5\linewidth}\centering
\img{.55}{deltaz_pu200.png}\\[.2em]\footnotesize $\Delta z$ between reconstructed vertices, PU200
\end{column}\begin{column}{.47\linewidth}\centering
\img{.55}{sigma_vs_mu.png}\\[.2em]\footnotesize $\svtx$ versus $\mu$
\end{column}\end{columns}
\par\vspace{.4em}
\takeaway{PU200: $\svtx=\sigPVF\pm\sigPVFErr$ mm for PV-Finder, $\sigAMVF\pm\sigAMVFErr$ mm for AMVF.}
\end{frame}

%----------------------------------------------------------------------
\begin{frame}{Hard-Scatter Vertex Reconstruction}
\begin{columns}[c]\begin{column}{.6\linewidth}\centering
\img{.66}{hs_reco_eff_vs_mu.png}
\end{column}\begin{column}{.37\linewidth}
\small A reconstructed vertex within $\hsWindowMM$ mm of the hard-scatter vertex.
\par\vspace{1.2em}\centering
\bigvalue{$\hsRecoPVF\%$}{PV-Finder at PU200}
\par\vspace{.8em}
{\Large $\hsRecoAMVF\%$}\\[.3em]\small AMVF
\end{column}\end{columns}
\end{frame}

\partdivider{III}{Track association and hard scatter}{After the GNN: categories, hard scatter and position}

%----------------------------------------------------------------------
\begin{frame}{Full Chain at PU200}
\begin{columns}[c]\begin{column}{.6\linewidth}\centering
\img{.66}{clean_yield_pu200.png}
\end{column}\begin{column}{.37\linewidth}
\small\textbf{Clean vertices / truth vertices}
\par\vspace{.6em}
Clean: $\geq70\%$ of the vertex tracks from one interaction.
\par\vspace{.8em}\centering
\begin{tabular}{lr}\toprule
PV-Finder + GNN & $\cleanTruthGNN\%$\\
Nearest $z_0$ & $\cleanTruthZnear\%$\\
$z_0$ within $3\sigma$ & $\cleanTruthZsig\%$\\
AMVF & $\cleanTruthAMVF\%$\\\bottomrule
\end{tabular}
\end{column}\end{columns}
\takeaway{The GNN more than doubles the clean-vertex yield of AMVF and of $z_0$-based assignment.}
\end{frame}

%----------------------------------------------------------------------
\begin{frame}{Full Chain across Pile-up}
\begin{columns}[c]\begin{column}{.65\linewidth}\centering
\img{.7}{clean_yield_vs_mu.png}
\end{column}\begin{column}{.32\linewidth}
\small Clean vertices / truth vertices versus $\mu$ for each stage of the chain and for AMVF.
\par\vspace{.8em}
Lower panel: ratio to AMVF.
\end{column}\end{columns}
\end{frame}

%----------------------------------------------------------------------
\begin{frame}{Hard-Scatter Vertex after Track Association}
\begin{columns}[c]\begin{column}{.47\linewidth}
\centering\small
\begin{tabular}{lrr}\toprule
PU200 & GNN & AMVF\\\midrule
Clean HS vertex & $\cleanHSGNN\%$ & $\cleanHSAMVF\%$\\
HS tracks on HS vertex & $\hsTrkGNN\%$ & $\hsTrkAMVF\%$\\
HS vertex purity & $\hsPurGNN\%$ & $\hsPurAMVF\%$\\
HS selected ($\sum p_T^2$) & $\hsSelGNN\%$ & $\hsSelAMVF\%$\\\bottomrule
\end{tabular}
\par\vspace{.8em}\raggedright\footnotesize
ATLAS definitions; tracks counted, not fit-weighted. HS vertex: most HS tracks. Selected: largest $\sum p_T^2$ and within $\hsWindowMM$ mm. AMVF: its stored primary vertex.
\end{column}\begin{column}{.5\linewidth}\centering
\img{.6}{hs_selection_eff_vs_mu.png}
\end{column}\end{columns}
\end{frame}

%----------------------------------------------------------------------
\begin{frame}{Hard-Scatter Longitudinal Resolution}
\begin{columns}[c]\begin{column}{.55\linewidth}\centering
\img{.66}{hs_resolution_vs_ntrk.png}
\end{column}\begin{column}{.42\linewidth}
\small
\textbf{Vertex fit on the GNN tracks}
\par\vspace{.5em}
Adaptive 3D fit with the beam-spot constraint, seeded at the PV-Finder peak.
\par\vspace{.5em}
ITk definition: iterative RMS of $z_{\rm reco}-z_{\rm HS}$, $\hsResNcommon$ events with the HS found by both.
\par\vspace{.8em}\centering
\bigvalue{$\hsResGNN\pm\hsResGNNErr~\mu$m}{PV-Finder + GNN}
\par\vspace{.5em}
{\Large $\hsResAMVF\pm\hsResAMVFErr~\mu$m}\\[.3em]\small AMVF
\end{column}\end{columns}
\end{frame}

\partdivider{IV}{Low pile-up and collision data}{Run 2 simulation, Run 2 and Run 3 data}

%----------------------------------------------------------------------
\begin{frame}{Low Pile-up: Simulation and Collision Data}
\begin{columns}[c]\begin{column}{.49\linewidth}\centering
\img{.5}{datamc_run2_resolution.png}\\[.2em]\footnotesize Run 2 data and simulation
\end{column}\begin{column}{.49\linewidth}\centering
\img{.5}{datamc_run3_resolution.png}\\[.2em]\footnotesize Run 3 data and simulation
\end{column}\end{columns}
\par\vspace{.3em}\centering\footnotesize
\begin{tabular}{lccc}\toprule
 & Simulation & Run 2 data & Run 3 data\\\midrule
PV-Finder $\svtx$ [mm] & $\sigMC$ & $\sigRunTwo$ & $\sigRunThree$\\\bottomrule
\end{tabular}
\par\vspace{.2em}
Simulation: efficiency $\effMC\%$, $\noRefMCtruth$ fake vertices per event; AMVF $\svtx=\sigMCAMVF$ mm.
\end{frame}

\partdivider{V}{Timing and deployment}{Athena integration, CPU and GPU}

%----------------------------------------------------------------------
\begin{frame}{Athena: Timing on CPU and GPU}
\begin{columns}[c]\begin{column}{.58\linewidth}\centering
\img{.62}{athena_cpu_gpu_amvf.png}
\end{column}\begin{column}{.39\linewidth}
\small
\textbf{Same events, tracks to vertices with track links}
\par\vspace{.6em}
Intel Xeon Silver 4216, NVIDIA T4.
\par\vspace{.6em}\centering
\begin{tabular}{lr}\toprule
At $\mu=200$ & ms/event\\\midrule
PV-Finder + GNN, GPU & $\athGPUms$\\
PV-Finder + GNN, CPU & $\athCPUms$\\
AMVF, CPU & $\sim\athAMVF$\\\bottomrule
\end{tabular}
\par\vspace{.6em}
\textcolor{pennred}{GPU chain $\sim\athGPUfactor\times$ faster than AMVF.}
\end{column}\end{columns}
\context{Implemented in Athena by Trey Warner. Tools and validation in the backup.}
\end{frame}

%----------------------------------------------------------------------
\begin{frame}{A Smaller Network for CPU}
\begin{columns}[t]\begin{column}{.52\linewidth}
\small\textbf{Same data and loss, narrower U-Net}
\par\vspace{.5em}\centering
\begin{tabular}{lrrr}\toprule
 & v6 & s5 & s8\\\midrule
U-Net channels & 280 & 112 & 80\\
CPU time [ms] & $\cpuVsixPad$ & $\cpuSFivePack$ & $\cpuSEightPack$\\
Speed-up & $1$ & $\spdSFive\times$ & $\spdSEight\times$\\
$\Delta$efficiency [pt] & & $\sFiveDEffA$ & $\sEightDEffA$\\
$\Delta$fake / event & & $\sFiveDFake$ & $\sEightDFake$\\\bottomrule
\end{tabular}
\par\vspace{.5em}\raggedright\footnotesize
ONNX Runtime, one thread, one PU200 event. Small models use packed inference: tracks only, no padding, identical output.
\end{column}\begin{column}{.45\linewidth}
\small\textbf{Projected Athena time at $\mu=200$, CPU}
\par\vspace{.5em}\centering
\begin{tabular}{lr}\toprule
Chain & ms/event\\\midrule
v6 + GNN & $\athCPUms$\\
s5 + GNN & $\sim\projSFive$\\
s8 + GNN & $\sim\projSEight$\\
AMVF & $\sim\athAMVF$\\\bottomrule
\end{tabular}
\par\vspace{.6em}\raggedright\footnotesize
Packed inference alone: $\spdVsixPack\times$ for v6 at no cost.
\end{column}\end{columns}
\takeaway{The small network brings the CPU chain to the speed of AMVF at $\mu=200$.}
\end{frame}

%----------------------------------------------------------------------
\begin{frame}{Summary}
\begin{itemize}\setlength\itemsep{.75em}
\item At PU200, PV-Finder finds $\effPVF\%$ of truth vertices, $+\effDiff$ points over AMVF, with fewer fakes.
\item Vertex--vertex resolution $\sigPVF$ mm against $\sigAMVF$ mm for AMVF.
\item With GNN track association: $\cleanTruthGNN\%$ clean vertices per truth vertex, against $\cleanTruthAMVF\%$ for AMVF.
\item Hard scatter: selected in $\hsSelGNN\%$ of events ($\hsSelAMVF\%$ AMVF); position resolution $\hsResGNN~\mu$m ($\hsResAMVF~\mu$m AMVF).
\item Validated on Run 2 and Run 3 collision data.
\item Running in Athena on CPU and GPU.
\end{itemize}
\par\vspace{.6em}
\textbf{Next:} freeze the HL-LHC density model and GNN, and complete the paper draft.
\end{frame}

%----------------------------------------------------------------------
\backupbegin

\begin{frame}{Backup: Efficiency Difference versus Matching Window}
\begin{columns}[c]\begin{column}{.6\linewidth}\centering
\img{.68}{window_scan_eff_diff.png}
\end{column}\begin{column}{.37\linewidth}
\small PV-Finder minus AMVF efficiency, same events, for windows from $0.1$ to $1$ mm.
\par\vspace{.8em}
PV-Finder is ahead at every window.
\end{column}\end{columns}
\end{frame}

\begin{frame}{Backup: Run 2 and HL-LHC Resolution}
\begin{columns}[c]\begin{column}{.52\linewidth}\centering
\img{.6}{sigma_summary.png}
\end{column}\begin{column}{.45\linewidth}
\small
\textbf{AMVF / PV-Finder $\svtx$:} $\ratioRunTwo$ in Run 2 simulation, $\sigRatioHL$ at PU200.
\par\vspace{.6em}
Track-based vertex precision: $\idealPURunTwo~\mu$m (Run 2), $\idealPUHL~\mu$m (HL-LHC).
\par\vspace{.6em}
\textbf{Training target through the same peak finding:}
\par\vspace{.3em}\centering\footnotesize
\begin{tabular}{lrr}\toprule
$\svtx$ [mm] & Run 2 & PU200\\\midrule
Target & $\tgtSigRunTwo$ & $\tgtSigHL$\\
PV-Finder & $\sigMCPVFr$ & $\sigPVF$\\
PV-Finder / target & $\pvfOverTgtRunTwo$ & $\pvfOverTgtHL$\\\bottomrule
\end{tabular}
\par\vspace{.4em}\raggedright
A $0.02$ mm grid moves the target from $\tgtSigHLwone$ to $\tgtSigHLfine$ mm.
\end{column}\end{columns}
\end{frame}

\begin{frame}{Backup: AMVF Hard-Scatter Selection}
\begin{columns}[c]\begin{column}{.55\linewidth}\centering
\img{.62}{amvf_hs_selection_vs_mu.png}
\end{column}\begin{column}{.42\linewidth}
\small
\textbf{AMVF stored primary vertex:} $\hsSelAMVF\%$ at PU200, as in the ITk performance paper.
\par\vspace{.6em}
Re-ranking the stored track lists by $\sum p_T^2$ gives $\hsSelAMVFrerank\%$: a few very high-$p_T$ tracks change the leading vertex.
\end{column}\end{columns}
\end{frame}

\begin{frame}{Backup: Hard-Scatter Tracks and Clean HS versus Pile-up}
\begin{columns}[c]\begin{column}{.49\linewidth}\centering
\img{.62}{hs_clean_fraction_vs_mu.png}
\end{column}\begin{column}{.49\linewidth}\centering
\img{.62}{hs_track_fraction_vs_mu.png}
\end{column}\end{columns}
\end{frame}

\begin{frame}{Backup: Low Pile-up Vertex Yield and the Reference}
\begin{columns}[c]\begin{column}{.52\linewidth}\centering
\img{.62}{datamc_run2_yield.png}
\end{column}\begin{column}{.45\linewidth}
\small
\textbf{PV-Finder peaks without a reference vertex, per event}
\par\vspace{.5em}\centering
\begin{tabular}{lr}\toprule
Simulation, truth reference & $\noRefMCtruth$\\
Simulation, AMVF reference & $\noRefMCamvf$\\
Run 2 data, AMVF reference & $\noRefRunTwo$\\
Run 3 data, AMVF reference & $\noRefRunThree$\\\bottomrule
\end{tabular}
\par\vspace{.6em}\raggedright
In data the reference is AMVF; simulation with the same reference reproduces the data.
\end{column}\end{columns}
\end{frame}

\begin{frame}{Backup: Athena Integration}
\relax
{\centering\resizebox{.94\linewidth}{!}{%
\begin{tikzpicture}[x=1cm,y=1cm,font=\small]
\foreach \x/\id/\txt in {0/trk/{xAOD tracks},3.4/sel/{Track selection\\+ tensor building},6.8/net/{ONNX inference\\+ peak finding},10.2/vtx/{xAOD vertices}}{
\node[pvbox,draw=mydarkblue!60,fill=mydarkblue!7,minimum width=2.9cm,minimum height=1.1cm] (\id) at (\x,0) {\txt};}
\draw[pvflow] (trk)--(sel);\draw[pvflow] (sel)--(net);\draw[pvflow] (net)--(vtx);
\node[pvbox,draw=pennred!70,fill=pennred!7,text width=5.0cm,minimum height=1.0cm] (gnn) at (6.8,-2.0) {GNN: build graph $\to$ score edges\\$\to$ select track links};
\draw[pvflow] (trk.south) -- (0,-.95) -| ([xshift=-12mm]gnn.north);
\draw[pvflow] (vtx.south) -- (10.2,-2.0) -- (gnn.east);
\node[pvbox,draw=mygreen!60!black,fill=mygreen!7,text width=3.1cm,minimum height=1.0cm] (links) at (0,-2.0) {Vertices with\\linked tracks};
\draw[pvflow,draw=mygreen!60!black] (gnn.west)--(links.east);
\end{tikzpicture}}\par}
\par\vspace{.8em}\small
\texttt{InDetNNPriVxFinderTool} and \texttt{InDetGNNTrackVertexAssociationTool}, with the same \texttt{IVertexFinder} interface as AMVF. ONNX Runtime on CPU or CUDA, set by configuration. Athena output agrees with the Python implementation to floating-point precision.
\par\vspace{.5em}{\scriptsize\color{black!55}\treysource.}
\end{frame}

\backupend
\end{document}
